% ISEG MASTER DISSERTATION
% MASTER'S IN MATHEMATICAL FINANCE
% AUTHOR: Petr Terletskiy
% SUPERVISOR: Prof. João Afonso Bastos
%
% Based on the ISEG template by Costa (2017),
% adapted by Domingos Guerreiro Seward and Luís F. Costa (2018).
%
% NOTE: Compile with pdfLaTeX (or XeLaTeX if fontspec is used)
%
% =============================================
% Preamble
% =============================================

\documentclass[12pt,a4paper,oneside]{article}

% Font
\usepackage{times}
\usepackage[T1]{fontenc}
\usepackage[utf8]{inputenc}

% Language
\usepackage[english]{babel}

% Appendix
\usepackage[toc,page]{appendix}

\makeatletter
\def\@biblabel#1{\hspace*{-\labelsep}}
\makeatother

% Bibliography
\usepackage{natbib}

% Margins (3cm all sides per ISEG guidelines)
\usepackage[top=3cm, bottom=3cm, left=3cm, right=3cm]{geometry}

% Change margins command
\def\changemargin#1#2{\list{}{\rightmargin#2\leftmargin#1}\item[]}
\let\endchangemargin=\endlist

% Paragraph indentation
\usepackage{indentfirst}

% 1.5 line spacing
\usepackage{setspace}
\onehalfspacing

% Paragraph spacing
\parskip=6pt

% Headers and footers
\usepackage{fancyhdr}
\pagestyle{fancy}
\fancyhf{}
\lhead{\footnotesize \textsc{Petr Terletskiy}}
\rhead{\footnotesize\textsc{BTC Direction Prediction Using KANs}}
\cfoot{\thepage}

% Section formatting
\usepackage{titlesec}
\titleformat{\chapter}[display]
  {\scshape\center}
  {}
  {0pt}
  {\thechapter.\ }
\usepackage{etoolbox}
\makeatletter
\patchcmd{\chapter}{\if@openright\cleardoublepage\else\clearpage\fi}{}{}{}
\makeatother
\titleformat{\section}
  {\scshape\center}{\thesection}{1em}{}
\titleformat{\subsection}
  {\itshape\center}{\thesubsection}{1em}{}
\titleformat{\subsubsection}
  {\itshape\center}{\thesubsubsection}{1em}{}

% Hyperlinks
\usepackage{hyperref}
\hypersetup{
    colorlinks=true,
    linkcolor=black,
    citecolor=black,
    urlcolor=blue
}

% Table of contents
\usepackage{tocloft}
\setlength{\cftbeforetoctitleskip}{0pt}
\setlength{\cftbeforeloftitleskip}{0pt}
\setlength{\cftbeforelottitleskip}{0pt}

% Tables
\renewcommand*\thetable{\Roman{table}}
\addto\captionsenglish{\renewcommand{\tablename}{\textsc{Table}}}
\addto\captionsenglish{\renewcommand{\figurename}{\textsc{Figure}}}

\usepackage{lipsum}
\usepackage{booktabs}
\usepackage{adjustbox}
\setlength\heavyrulewidth{0.3ex}
\usepackage{tabularx}
\usepackage{array}
\usepackage{threeparttable}
\usepackage{multirow}
\usepackage{longtable}

% Quotes
\usepackage{csquotes}

% Mathematics
\usepackage{amsmath}
\usepackage{amsfonts}
\usepackage{amssymb}
\usepackage{bm}

% Graphics
\usepackage{graphicx}
\usepackage{float}
\usepackage{subcaption}

% Code listings (for algorithm pseudocode)
\usepackage{algorithm}
\usepackage{algpseudocode}

% Glossary / Acronyms
\usepackage[automake]{glossaries}
\makeglossaries

% ── Acronyms ──────────────────────────────────────────────────────────
% AFML \& Methodology
\newacronym{afml}{AFML}{Advances in Financial Machine Learning}
\newacronym{tbl}{TBL}{Triple Barrier Labeling}
\newacronym{cusum}{CUSUM}{Cumulative Sum}
\newacronym{ffd}{FFD}{Fixed-width Window Fractional Differentiation}
\newacronym{cpcv}{CPCV}{Combinatorial Purged Cross-Validation}
\newacronym{dsr}{DSR}{Deflated Sharpe Ratio}
\newacronym{pbo}{PBO}{Probability of Backtest Overfitting}
\newacronym{mda}{MDA}{Mean Decrease Accuracy}

% Models
\newacronym{kan}{KAN}{Kolmogorov-Arnold Network}
\newacronym{lstm}{LSTM}{Long Short-Term Memory}
\newacronym{rf}{RF}{Random Forest}
\newacronym{xgb}{XGBoost}{Extreme Gradient Boosting}
\newacronym{lr}{LR}{Logistic Regression}
\newacronym{ar}{AR}{Autoregressive}
\newacronym{mlp}{MLP}{Multi-Layer Perceptron}

% Technical
\newacronym{btc}{BTC}{Bitcoin}
\newacronym{ohlcv}{OHLCV}{Open-High-Low-Close-Volume}
\newacronym{ta}{TA}{Technical Analysis}
\newacronym{ewma}{EWMA}{Exponentially Weighted Moving Average}
\newacronym{adf}{ADF}{Augmented Dickey-Fuller}
\newacronym{rsi}{RSI}{Relative Strength Index}
\newacronym{macd}{MACD}{Moving Average Convergence Divergence}
\newacronym{atr}{ATR}{Average True Range}
\newacronym{obv}{OBV}{On-Balance Volume}
\newacronym{mfi}{MFI}{Money Flow Index}
\newacronym{cci}{CCI}{Commodity Channel Index}
\newacronym{vix}{VIX}{CBOE Volatility Index}
\newacronym{mvrv}{MVRV}{Market Value to Realized Value}
\newacronym{sadf}{SADF}{Supremum Augmented Dickey-Fuller}

% Financial
\newacronym{mfw}{MFW}{Master's Final Work}
\newacronym{jel}{JEL}{Journal of Economic Literature}
\newacronym{emh}{EMH}{Efficient Market Hypothesis}


% =============================================
% Begin Document
% =============================================

\begin{document}

% ── 1. Cover Page ─────────────────────────────────────────────────────
\begin{titlepage}
    \begin{flushleft}
        \includegraphics[width=5cm]{ISEG_logo.png}
    \end{flushleft}
    
    \vspace{1.5cm}
    
    \begin{center}
        {\Large \textsc{Master}}\\[0.5cm]
        {\Large \textsc{MATHEMATICAL FINANCE}}\\[2cm]
        
        {\Large \textsc{Master's Final Work}}\\[0.5cm]
        {\Large \textsc{PROJECT}}\\[3cm]
        
        {\LARGE \textbf{BITCOIN DIRECTION PREDICTION USING KOLMOGOROV-ARNOLD NETWORKS WITHIN THE AFML FRAMEWORK}}\\[3cm]
        
        {\Large \textsc{Petr Terletskiy}}\\[1.5cm]
                
        \vfill
        
        {\large \textsc{May 2026}}
    \end{center}
\end{titlepage}

% ── 2. Title Page (Inside Cover) ──────────────────────────────────────
\begin{titlepage}
    \begin{flushleft}
        \includegraphics[width=5cm]{ISEG_logo.png}
    \end{flushleft}
    
    \vspace{1.5cm}
    
    \begin{center}
        {\Large \textsc{Master}}\\[0.5cm]
        {\Large \textsc{MATHEMATICAL FINANCE}}\\[1.5cm]
        
        {\Large \textsc{Master's Final Work}}\\[0.5cm]
        {\Large \textsc{PROJECT}}\\[2cm]
        
        {\LARGE \textbf{BITCOIN DIRECTION PREDICTION USING KOLMOGOROV-ARNOLD NETWORKS WITHIN THE AFML FRAMEWORK}}\\[2cm]
        
        {\Large \textsc{Petr Terletskiy}}\\[2cm]
        
        {\large \textsc{Supervision:}}\\[0.2cm]
        {\Large \textsc{Prof. João Afonso Bastos}}\\[2cm]
        
        \vfill
        
        {\large \textsc{May 2026}}
    \end{center}
\end{titlepage}

% ── 3. Blank Page ─────────────────────────────────────────────────────
\newpage
\thispagestyle{empty}
\mbox{}

% ── 4. Dedication ─────────────────────────────────────────────────────
\newpage
\thispagestyle{empty}
\vspace*{\fill}
\begin{flushright}
\textit{%
To my family and God.
}
\end{flushright}

% ── 5. Erratum (commented out unless needed for final post-viva) ──────
%\newpage
%\thispagestyle{plain}
%\addcontentsline{toc}{section}{Erratum}
%\section*{Erratum}

% ── 6. Glossary ───────────────────────────────────────────────────────
\newpage
\thispagestyle{plain}
\addcontentsline{toc}{section}{Glossary}

\setcounter{page}{1}
\pagenumbering{roman}

\glsaddall
\printglossaries


% ── 7. Abstract ───────────────────────────────────────────────────────
\newpage
\thispagestyle{plain}
\addcontentsline{toc}{section}{Abstract, Keywords, and JEL Codes}

\section*{Abstract, Keywords, and \gls{jel} Codes}

% TODO: Write abstract (300 words max).
%
% Topics to cover, mapping the paper_outline narrative:
%   - Frame around Q1 to Q4 / C1 to C4 from the outline.
%   - Headline numbers from the locked run (v6, 5 seeds across all six models,
%     840 prediction entries, 1,168 aligned events, base rate 0.5685):
%       * No model achieves DSR >= 0.95 (top DSR among trained models =
%         0.2470, KAN).
%       * Buy-and-hold posts a higher median Sharpe (2.2722) than every
%         trained model; DSR is not defined for buy-and-hold since it
%         does not enter the Optuna selection pool.
%       * PBO = 0.657 lands in the adversarial regime (23 of 35 IS/OOS
%         partitions); the IS-best is the OOS-loser in nearly two of three
%         partitions. 6 of 7 bootstrap 95\% CIs cross zero; only
%         buy-and-hold's CI (1.8903, 2.4685) stays strictly positive.
%         The previous run's KAN strict-positive CI (0.0971, 1.2413)
%         does NOT reproduce under uniform 5-seed evaluation; KAN's CI
%         is now (-0.4677, 1.1607). Leave-one-out PBO shows Random
%         Forest is the PBO-anchor (Delta +0.000), while every other
%         model's exclusion reduces PBO; Logistic Regression is the
%         largest destabiliser (Delta -0.286 to 0.371).
%       * DeLong: 2 of 15 pairwise AUC tests reach alpha=0.05 (AR
%         Logistic vs XGBoost p=0.0427, Random Forest vs XGBoost
%         p=0.0318), both with XGBoost as the lower-AUC half. The
%         first locked configuration with any DeLong-significant pair.
%       * Symbolic extraction succeeds: 100\% symbolification rate (15
%         of 15 edges, no skips), three surviving features
%         (eth\_btc\_ratio, log\_returns\_lag6, natgas\_ret\_30),
%         post-symbolic accuracy 51.85\% (a 5.18 pp gain over
%         pre-symbolic 46.67\%, but both numbers sit at or below the 50\%
%         baseline). The symbolic-extraction outputs are identical to
%         the previous 3-seed run (PyKAN retrained with fixed seed on
%         same training fold). The Phase-1 dynamic gate (max(0.53,
%         base\_rate + 0.01) = 0.5693) fails at val\_acc 0.4815 and the
%         pre-prune gate also fails; the locked-run PyKAN model has not
%         extracted directional signal on the deployment fold. The
%         formula uses six clean primitives (sin, cos, tanh, x\^2, x\^3,
%         x\^4); the previous configuration's Abs(.) non-
%         differentiability flag is resolved and all surviving features
%         have finite gradients at the dataset mean. A lag feature
%         survives extraction for the second consecutive locked
%         configuration (log\_returns\_lag1 -> log\_returns\_lag6).
%   - Position the symbolic contribution as the first closed-form
%     classification formula extracted from a CPCV-trained KAN under
%     the full AFML evaluation.
This dissertation evaluates whether Kolmogorov-Arnold Networks can predict Bitcoin daily price direction and extract interpretable symbolic formulas, benchmarked against five models within López de Prado's Advances in Financial Machine Learning framework. Using Combinatorial Purged Cross-Validation with the Deflated Sharpe Ratio and Probability of Backtest Overfitting, the study finds that no model achieves statistically significant predictive performance, consistent with market efficiency. The core contribution is KAN symbolic expression extraction, deriving closed-form decision functions from the trained network.

\vspace{1em}

\textsc{Keywords}: Bitcoin; Kolmogorov-Arnold Networks; Symbolic Extraction; Triple Barrier Labeling; Combinatorial Purged Cross-Validation; Financial Machine Learning.

\vspace{1em}

\textsc{\gls{jel} Codes}: C45; C53; C58; G11; G17.


% ── 8. Table of Contents ──────────────────────────────────────────────
\newpage
\thispagestyle{plain}
\addcontentsline{toc}{section}{Table of Contents}
\renewcommand*\contentsname{\hfill\normalfont\scshape\normalsize Table of Contents \hfill}
\tableofcontents

% ── 9. List of Figures ────────────────────────────────────────────────
\newpage
\thispagestyle{plain}
\addcontentsline{toc}{section}{List of Figures}
\renewcommand*\listfigurename{\hfill\normalfont\scshape\normalsize List of Figures \hfill}
\listoffigures

% ── 10. List of Tables ────────────────────────────────────────────────
\newpage
\thispagestyle{plain}
\addcontentsline{toc}{section}{List of Tables}
\renewcommand*\listtablename{\hfill\normalfont\scshape\normalsize List of Tables \hfill}
\listoftables

% ── 11. Preface (optional) ────────────────────────────────────────────
%\newpage
%\thispagestyle{plain}
%\addcontentsline{toc}{section}{Preface}
%\section*{Preface}


% ── 12. Acknowledgements ──────────────────────────────────────────────
\newpage
\thispagestyle{plain}
\addcontentsline{toc}{section}{Acknowledgements}
\section*{Acknowledgements}

% TODO: Write acknowledgements.
First, I wish to thank Professor João Afonso Bastos for his encouragement and guidance. I am also thankful to my family for their patience and support while I pursued this project.


% =============================================
% TEXTUAL PART BEGINS
% =============================================
\newpage
\thispagestyle{plain}

\setcounter{page}{1}
\pagenumbering{arabic}

\begin{center}
\textsc{BITCOIN DIRECTION PREDICTION USING KOLMOGOROV-ARNOLD NETWORKS WITHIN THE AFML FRAMEWORK}\\

\vspace{0.5cm}
\textsc{BY PETR TERLETSKIY}\\
\end{center}


% ==================================================================
% CHAPTER 1: INTRODUCTION
% ==================================================================
\section{Introduction}\label{sec:introduction}

% Target: max 20\% of textual part (~7 pages for 35-page MFW).
% Cochrane: "Start with your central contribution."
% Thatcher R1: "First 1-2 sentences state a clear, field-relevant claim."
% No subsection number for the opening paragraphs.
%
% ── OPENING PARAGRAPHS (before any subsection) ───────────────────
%
% Three sentences, in order:
%   - This thesis applies KANs to BTC daily direction prediction within the
%     full AFML framework (CUSUM + TBL + FFD + sample weights + CPCV +
%     DSR + PBO + DeLong).
%   - It benchmarks KAN against five models (AR Logistic, LR, RF, XGBoost,
%     LSTM) under CPCV with N=8, k=2, 28 splits, and 7 backtest paths;
%     the locked configuration uses 5 seeds across all six models,
%     producing 840 prediction entries.
%   - The novel deliverable is a closed-form symbolic decision function
%     extracted from the trained KAN.
%
% Telegraph the headline findings immediately (T-R1, T-R3) using the
% locked-run v6 numbers:
%   - No model achieves DSR >= 0.95 (top DSR among trained models =
%     0.2470, KAN). Buy-and-hold posts a higher median Sharpe (2.2722)
%     than every trained model but is excluded from DSR since it does
%     not enter the selection pool.
%   - PBO = 0.657 sits in the adversarial regime (23 of 35 IS/OOS
%     partitions); the IS-best is the OOS-loser in nearly two of three
%     partitions. 6 of 7 bootstrap 95\% CIs cross zero; only
%     buy-and-hold's CI (1.8903, 2.4685) stays strictly positive.
%     The previous run's KAN strict-positive CI (0.0971, 1.2413) does
%     NOT reproduce under uniform 5-seed evaluation; KAN's CI is now
%     (-0.4677, 1.1607). Leave-one-out PBO shows Random Forest is the
%     PBO-anchor (Delta +0.000), while every other model's exclusion
%     reduces PBO (logistic -0.286, ar\_logistic -0.200, lstm -0.200,
%     xgboost -0.114, kan -0.114). Random Forest sits low enough in
%     the IS ranking to be PBO-neutral; Logistic Regression is the
%     largest destabiliser.
%   - DeLong: 2 of 15 pairwise AUC tests reach alpha=0.05 (AR Logistic
%     vs XGBoost p=0.0427, Random Forest vs XGBoost p=0.0318), both
%     with XGBoost as the lower-AUC half. The first locked
%     configuration with any DeLong-significant pair; the 5-seed
%     averaging reduces the noise floor enough for the largest AUC
%     gaps to clear.
%   - The symbolic extraction pipeline succeeds with 100\% symbolification
%     rate (15 of 15 edges, no skips); three features survive
%     (eth\_btc\_ratio, log\_returns\_lag6, natgas\_ret\_30);
%     post-symbolic accuracy 51.85\% (a 5.18 pp gain over pre-symbolic
%     46.67\%, but both numbers sit at or below the 50\% baseline).
%     The symbolic-extraction outputs are identical to the previous
%     3-seed run (PyKAN retrained with fixed seed on same training
%     fold). The formula uses six clean primitives (sin, cos, tanh,
%     x\^2, x\^3, x\^4); the previous configuration's Abs(.)
%     non-differentiability flag is resolved. A lag feature survives
%     extraction for the second consecutive locked configuration
%     (log\_returns\_lag6 displaces the previous run's
%     log\_returns\_lag1). The Phase-1 dynamic gate (max(0.53,
%     base\_rate + 0.01) = 0.5693) fails at val\_acc 0.4815; the
%     pre-prune gate also fails, and recovery in the LBFGS phases
%     does not push validation accuracy above the 50\% baseline.
%
% KAN reference here is one sentence only. No theorem, no splines, no
% architecture details. Save those for Section 2.1.


\subsection{Research Context and Motivation}\label{subsec:context}

Bitcoin was presented to the world on October 31, 2008 by Satoshi Nakamoto. Since its first transaction on January 12, 2009, Bitcoin has produced one of the highest cumulative returns of any major asset, and it gave birth to the cryptocurrency asset class. It is today the largest cryptocurrency by market capitalization. Unlike traditional asset classes, Bitcoin trades 365 days a year, 24/7 on a transparent and publicly auditable blockchain, exposing on-chain activity that has no equivalent in equity or fixed-income markets.
% The narrative hook (T-R2): "despite A, we observe B".
% A: crypto-ML literature reports 85\% to 95\% daily-direction accuracy
%    using deep learning architectures.
% B: no model in this thesis survives the AFML statistical-correction layer.
%
% ── Why BTC ────────────────────────────────────────────────────────
%
% One paragraph, factual justification only:
%   - Largest cryptocurrency by market capitalization.
%   - Highest cumulative return of any asset class since 2009.
%   - Trades 24/7 on a transparent, publicly auditable blockchain;
%     on-chain features have no equivalent in traditional finance.
%   - Daily OHLCV data freely available from 2014, providing roughly
%     4,200 daily observations.
%
% ── Why KANs ───────────────────────────────────────────────────────
%
% One paragraph:
%   - Liu et al. (2024) introduced an MLP alternative whose edges carry
%     learnable B-splines.
%   - Those splines can be distilled into closed-form symbolic functions.
%   - No prior work applies KANs to BTC daily direction or extracts
%     symbolic formulas from a classification KAN under AFML evaluation.


\subsection{Problems and Research Questions}\label{subsec:objectives}

% Q1 to Q4 mapped one-to-one onto C1 to C4 (Thatcher R6: parallel logic
% across Introduction, Methodology, Results, Discussion).
% Two short paragraphs each: state the problem in the field, then state
% the question this thesis asks.
%
%
% ── Q1. Predictability under leakage-free evaluation ──────────────
%
% Problem:
%   - Crypto-ML papers report 85\% to 95\% daily-direction accuracy using
%     fixed-horizon labels with overlapping spans and naive train-test
%     splits with no purging or embargo.
% Question:
%   - Does BTC daily direction remain predictable once labels stop leaking
%     the future and overlapping observations are downweighted?
%   - Does any positive Sharpe survive the AFML statistical-correction
%     layer (DSR, PBO, DeLong)?
%
%
% ── Q2. Which feature families carry signal ───────────────────────
%
% Problem:
%   - Most crypto-ML papers use one or two feature families (typically
%     TA or LSTM-on-prices).
%   - No consensus on whether macro, crypto-macro, or on-chain features
%     add information beyond price-derived features.
% Question:
%   - Among technical, statistical, macroeconomic, crypto-macro, and
%     on-chain features, which families survive multi-model permutation-
%     importance selection across CPCV folds?
%   - Are on-chain features the "free lunch" the literature suggests?
%
%
% ── Q3. KAN versus standard model families ────────────────────────
%
% Problem:
%   - KANs have been applied to VIX forecasting and stock prediction,
%     but never to BTC direction.
%   - Their position relative to econometric, linear, tree-based, and
%     deep-learning baselines under a uniform protocol is unknown.
% Question:
%   - Where does a KAN sit relative to AR Logistic, Logistic Regression,
%     Random Forest, XGBoost, and LSTM under identical CPCV splits,
%     identical features, identical sample weights, and identical metrics?
%
%
% ── Q4. Closed-form formula extraction ────────────────────────────
%
% Problem:
%   - Existing KAN symbolic-extraction work targets regression problems
%     on relatively predictable series (VIX, stock prices).
%   - No work has extracted a closed-form classification formula from
%     a CPCV-trained KAN in a weak-signal regime.
% Question:
%   - Can a human-readable mathematical expression for P(Up) be extracted
%     from a CPCV-trained KAN, while preserving most of the predictive
%     accuracy of the trained network?


\subsection{Contributions}\label{subsec:contributions}

% Map one-to-one onto Q1 to Q4. Keep prose tight; methodology details
% live in Section 3.
%
%
% ── C1 -> Q1. Honest evaluation under full AFML ───────────────────
%
%   - End-to-end pipeline applying the complete AFML stack: CUSUM event
%     sampling, triple-barrier labels, sample weights (uniqueness x
%     return attribution x time decay), fractional differencing, CPCV
%     with purging and embargo, plus DSR, PBO, and DeLong corrections.
%
%
% ── C2 -> Q2. A 73-feature universe across four families ──────────
%
%   - 25 technical, 9 statistical (AFML Part 4), 29 external (20 macro,
%     1 crypto-macro, 8 on-chain from CoinMetrics), plus 10 autoregressive
%     lag features.
%   - All 73 features compete in multi-model MDA selection.
%
%
% ── C3 -> Q3. Six-model apples-to-apples benchmark ────────────────
%
%   - Identical CPCV splits, identical features, identical sample weights,
%     identical metrics across all six models.
%   - Pairwise DeLong AUC tests determine which differences are
%     statistically real.
%
%
% ── C4 -> Q4. Closed-form symbolic formula from KAN ───────────────
%
%   - A human-readable expression for P(Up), pruned and substituted from
%     the trained KAN via PyKAN's symbolic-extraction pipeline.
%   - The novel contribution that distinguishes this thesis from a standard
%     model-benchmarking study.


\subsection{Document Structure}\label{subsec:structure}

% Cochrane suggests skipping a roadmap; ISEG conventions favour a brief
% one. Keep it minimal: name Sections 2 to 7 in two lines max.



% ==================================================================
% CHAPTER 2: LITERATURE REVIEW
% ==================================================================
\newpage
\section{Literature Review}\label{sec:literature}

% Cochrane: position against the 2 to 3 closest papers per stream.
% Three streams: KANs, BTC ML prediction, AFML methodology.
% Each stream ends with a gap statement.
%
% Placement notes from the outline:
%   - HARd to Beat (Chassot \& Audrino 2026) goes in 2.3.9 (AFML),
%     NOT in BTC ML. The paper is methodological; its claim
%     ("fitting scheme matters more than model choice") supports
%     this thesis's DSR=0 narrative directly.
%   - Kolmogorov-Arnold theorem and splines go in 2.1 only. NOT in
%     Introduction (no math before main result). NOT in Methodology
%     (Methodology assumes the reader has read 2.1).


\subsection{Kolmogorov-Arnold Networks}\label{subsec:lit_kan}

% Section opening claim (T-R3): "KANs are the only neural architecture
% whose learned activation functions can be distilled into closed-form
% symbolic expressions." State this first, then defend.


\subsubsection{Mathematical foundation}\label{subsubsec:kan_math}

% Topics:
%   - Kolmogorov-Arnold representation theorem (Kolmogorov 1957;
%     Arnold 1958): any continuous multivariate function f(x_1, ..., x_n)
%     decomposes as a finite sum of continuous univariate functions.
%   - The theorem was historically considered impractical because the
%     inner functions can be highly non-smooth.
%   - Liu et al. (2024) made it practical by parameterizing those
%     functions with learnable B-splines.
%
% Cite: \cite{kolmogorov_1957}, \cite{arnold_1958}.


\subsubsection{How KANs work}\label{subsubsec:kan_arch}

% Topics:
%   - MLPs put fixed activation functions on nodes; KANs put learnable
%     activation functions on edges.
%   - Each edge is a B-spline g(x) = sum a_k * B_k(x) of order k over
%     G grid intervals.
%   - Architecture notation: [n_input, n_hidden, n_output].
%   - Training: Adam followed by L-BFGS, with L1-on-activations and
%     entropy sparsity regularisation. Adam escapes local minima,
%     L-BFGS exploits curvature once Adam has placed the model in a
%     good basin.
%   - Training best practices on small datasets (Noorizadegan 2026):
%     coarser grids generalise better, finer grids overfit; tanh input
%     normalisation to [-1, 1] matches the default grid\_range; L1 +
%     entropy regularisation aids pruning and symbolification. The
%     practical methodology choices for this thesis (single hidden
%     layer, grid in {3, 5}, etc.) live in Section 3.7.
%   - Key property exploited in this thesis: after training, each edge's
%     spline can be inspected, pruned, or replaced with a symbolic
%     function (sin, exp, x^2, tanh), yielding an interpretable
%     closed-form expression.
%
% Cite: \cite{liu_kan_2024}, \cite{noorizadegan_2026}.


\subsubsection{KAN variants}\label{subsubsec:kan_variants}

% Topics:
%   - KAN 2.0 / MultKAN (Liu et al. 2024b) introduces multiplication
%     nodes alongside addition. Enables discovery of multiplicative
%     interactions (e.g., RSI x Stoch\_K) without log-exp decomposition.
%     The same symbolic-extraction pipeline applies. MultKAN is not used
%     in the main results of this thesis and is cited as future-work
%     justification (Section 7).
%   - TKAN (Genet and Inzirillo 2024) is the canonical recurrent-KAN
%     variant: replaces LSTM gates with KAN edges. Cited as
%     representative of the architectural-variant literature this thesis
%     chose not to pursue, since sequence variants are incompatible with
%     PyKAN's symbolic-extraction pipeline.
%   - Comprehensive-review anchor. The Yamak et al. (2025) review of
%     KAN time-series variants covers TKAN, ChebyKAN, and roughly 40
%     further variants under a single survey, avoiding the need to cite
%     each variant individually. The review notes that symbolic
%     distillation in financial time series is underexplored, which is
%     the methodological gap this thesis addresses.
%
% Cite: \cite{liu_kan2_2024}, \cite{genet_inzirillo_2024},
% \cite{yamak_et_al_2025}.


\subsubsection{KANs for financial time series}\label{subsubsec:kan_finance}

% Three papers bracket what this thesis does and does not do.
%
% ── VIX KAN paper (Cho, Lee, and Kim 2025) ────────────────────────
%
% Topics:
%   - Applies KAN to VIX forecasting (regression).
%   - Introduces Algorithm 1: Train (Adam -> grid extend -> LBFGS) ->
%     Prune -> Symbolify -> Affine fine-tune.
%   - Direct methodological ancestor of this thesis's symbolic extraction
%     pipeline.
%   - Key difference: VIX is mean-reverting and predictable; BTC daily
%     direction is binary classification in a weak-signal regime, which
%     makes symbolic extraction substantially harder.
%
% Cite: \cite{cho_lee_kim_2025}.
%
% ── KASPER (Oad et al. 2025) ──────────────────────────────────────
%
% Topics:
%   - KAN combined with regime detection (Gumbel-Softmax mechanism)
%     for stock prediction.
%   - Reports R^2 = 0.89 and Sharpe = 12.02 on individual stocks.
%   - Closest related work in terms of KAN + finance + interpretability.
%   - Differences from this thesis: regression not classification; regime
%     detection layer absent here; Shapley-based rule extraction vs.
%     direct formula; no AFML evaluation.
%
% Cite: \cite{oad_kasper_2025}.
%
% ── DecoKAN (Gao et al. 2025) ─────────────────────────────────────
%
% Topics:
%   - DecoKAN applies KAN with an interpretable decomposition (trend,
%     seasonality, residual) to cryptocurrency forecasting. Closest to
%     this thesis on the asset-class axis (crypto), but uses a regression
%     target with naive train-test splits and no AFML correction layer.
%   - Together with VIX KAN and KASPER, this reinforces the gap
%     statement: KAN-on-crypto exists, KAN + symbolic extraction exists,
%     KAN + AFML + symbolic distillation under leakage-free evaluation
%     does not.
%
% Cite: \cite{gao_decokan_2025}.


\subsubsection{Gap statement}\label{subsubsec:kan_gap}

% Topics:
%   - No prior work applies KANs to BTC daily direction.
%   - No prior work extracts symbolic formulas from a classification KAN.
%   - No prior work combines KAN symbolic extraction with AFML evaluation.


\subsection{Bitcoin Price Prediction with Machine Learning}\label{subsec:lit_btc}

% Section opening claim (T-R3):
% "Despite extensive ML coverage of BTC prediction, the field's reported
% accuracy numbers do not survive leakage-free evaluation."


\subsubsection{Overview of approaches}\label{subsubsec:btc_overview}

% Topics:
%   - Three families: statistical/econometric (ARIMA, GARCH, VAR);
%     classical ML (RF, XGBoost, SVM, LR); deep learning (LSTM, GRU,
%     CNN, Transformer).
%   - Two framings: price-level regression or direction classification
%     (this thesis).


\subsubsection{Technical analysis and ML for BTC}\label{subsubsec:btc_ta}

% Topics:
%   - Confluence of TA and ML for BTC prediction (Máté et al. 2024).
%   - Identifies ROC as a top predictor (justifies inclusion of roc\_14
%     in this thesis's feature set).
%   - Demonstrates multi-indicator strategies outperform single-indicator.
%   - Limitation: standard split, no purging, no overlap accounting.
%
% Cite: \cite{mate_confluence_2024}.


\subsubsection{Deep learning for BTC direction}\label{subsubsec:btc_dl}

% Topics:
%   - Same problem as this thesis (BTC daily direction with DL models).
%     Reports competitive LSTM accuracy.
%   - Limitation: fixed-horizon labels, no CUSUM, no sample weights, no
%     purging or embargo.
%   - Comparing conclusions illustrates how methodology shapes reported
%     performance.
%
% Cite: \cite{omole_enke_2024}.


\subsubsection{Broader crypto-DL landscape and conservative prediction}\label{subsubsec:btc_dl_landscape}

% Topics:
%   - Broader landscape. Catch-all citations for the broader crypto-DL
%     space, to avoid citing 20+ individual papers. Documents LSTM, GRU,
%     and attention-based dominance. Documents recurring issues: small
%     datasets, overfitting, no transaction costs, no risk-adjusted
%     metrics, inconsistent evaluation.
%   - Conservative prediction in noisy financial regimes. Nabar and
%     Shroff (2023) argue that in noisy financial settings the right
%     behaviour is selective abstention rather than always-predict. A
%     model that abstains when confidence is low can outperform one
%     that predicts on every observation. Directly motivates the
%     bet-sizing threshold mechanism in Section 3.10.1: predictions
%     with p ~ 0.50 produce bet ~ 0, which is structural abstention
%     rather than a separate model.
%
% Cite: \cite{bourday_crypto_dl_2024}, \cite{wu_crypto_dl_review_2024},
% \cite{nabar_shroff_2023}.


\subsubsection{Methodological weaknesses in the literature}\label{subsubsec:btc_weaknesses}

% Bridge to Section 2.3. Synthesise the common problems across the papers
% reviewed above:
%   - Naive train-test splits with overlapping labels leaking future
%     information.
%   - Fixed-horizon labelling ignores stop-loss and take-profit dynamics.
%   - No sample weighting; redundant overlapping events count equally.
%   - Accuracy as primary metric ignores class balance and economic
%     significance.
%   - No correction for multiple testing, so reported Sharpes are
%     inflated.
%
% These are the failure modes that the AFML framework was designed to
% address (next subsection).


\subsection{The AFML Methodology}\label{subsec:lit_afml}

% Section opening claim (T-R3): "AFML treats financial ML as a
% methodology problem first and a model problem second."


\subsubsection{Overview and motivation}\label{subsubsec:afml_overview}

% Topics:
%   - López de Prado (2018) identifies systematic flaws in financial ML:
%     backtest overfitting, label-overlap leakage, naive CV, inflated
%     Sharpes from multiple testing.
%   - Proposes a complete framework addressing each flaw. Asset-agnostic
%     in principle.
%
% Cite: \cite{lopez_de_prado_2018}.


\subsubsection{Event-driven sampling, triple barrier labelling, and sample weights}\label{subsubsec:afml_cusum_tbl_weights}

% This subsection covers the three pre-CPCV AFML transformations that
% turn raw bars into (X, y, w, t1) tuples: event sampling, labelling,
% and sample weighting.
%
% ── CUSUM event filter ────────────────────────────────────────────
%
% Topics:
%   - Symmetric CUSUM (AFML Snippet 2.4) identifies structurally
%     meaningful events.
%   - Reduces the dataset to informationally rich observations,
%     discarding low-activity periods.
%   - Threshold typically calibrated as a multiple of mean daily
%     volatility.
%
% ── Triple barrier labelling ──────────────────────────────────────
%
% Topics:
%   - Three barriers per event: upper (take-profit), lower (stop-loss),
%     vertical (time limit).
%   - Label = first barrier touched. Carries t1 (resolution timestamp)
%     for downstream purging.
%   - Reflects real trading outcomes, not arbitrary fixed-horizon
%     snapshots.
%   - Empirical validation outside crypto: Kang and Kim (2025) apply TBL
%     to Korean equities and report that TBL on raw OHLCV inputs
%     outperforms fixed-horizon labels for short-term direction
%     prediction; the TBL machinery transfers across asset classes.
%
% ── Sample weights ────────────────────────────────────────────────
%
% Topics:
%   - Concurrent labels -> uniqueness -> return-attribution weights ->
%     time decay.
%   - Corrects for redundancy in overlapping labels.
%   - Passed to classifiers via sample\_weight (sklearn) or per-sample
%     loss multiplier (PyTorch).
%
% Cite: \cite{lopez_de_prado_2018}, Ch. 2 to 4; \cite{kang_kim_2025}.


\subsubsection{Fractional differentiation}\label{subsubsec:afml_ffd}

% Topics:
%   - Integer differentiation achieves stationarity but destroys long
%     memory; price levels preserve memory but are non-stationary.
%   - FFD finds the minimum d* in (0, 1) that achieves stationarity (ADF
%     test) while preserving maximum memory.
%   - Fixed-width window for computational efficiency.
%   - d* estimated per fold on training data only, as a per-fold
%     preprocessing step inside the CPCV loop (Section 3.6.1).
%
% Cite: \cite{lopez_de_prado_2018}, Ch. 5.


\subsubsection{Combinatorial Purged Cross-Validation}\label{subsubsec:afml_cpcv}

% Topics:
%   - Standard k-fold is invalid: overlapping labels leak across folds.
%   - CPCV: N contiguous groups, k test groups -> C(N, k) splits,
%     phi(N, k) backtest paths.
%   - Purging removes training observations whose labels overlap with the
%     test period (three sufficient conditions).
%   - Embargo removes a buffer after each test group to prevent
%     serial-correlation leakage.
%   - Path matrix enables distribution-level Sharpe analysis.
%
% Cite: \cite{lopez_de_prado_2018}, Ch. 7 and 12.


\subsubsection{Deflated Sharpe Ratio and Probability of Backtest Overfitting}\label{subsubsec:afml_dsr_pbo}

% Topics:
%   - DSR: corrects observed Sharpe for selection bias across n\_trials
%     models, also corrects for non-normal returns (skew, kurtosis).
%     DSR > 0.95 means significant after multiple-testing correction.
%   - PBO: combinatorially symmetric CV. Split paths into IS/OOS halves;
%     check if IS-best underperforms OOS median. PBO ~= 0 means reliable
%     selection; PBO > 0.5 means adversarial.
%
% Cite: \cite{lopez_de_prado_2018}, Ch. 11 and 14.


\subsubsection{AFML in practice: crypto applications and the fitting-scheme finding}\label{subsubsec:afml_crypto_hard}

% Two threads bracket how the AFML toolkit gets used in practice and
% why methodology, not model choice, drives outcomes.
%
% ── AFML applied to crypto ────────────────────────────────────────
%
% Topics:
%   - Slepaczuk and Bieganowski (2024) combine FFD and TBL with supervised
%     autoencoders on BTC, ETH, and LTC. Walk-forward d* estimation.
%     Validates that FFD + TBL transfer to crypto with improved
%     risk-adjusted returns.
%   - Fu et al. (2024) extend the AFML toolkit further: a genetic-
%     algorithm-driven TBL that searches over barrier widths and time
%     horizons jointly with model hyperparameters, applied to a pair-
%     trading strategy on cryptocurrencies. Demonstrates that TBL
%     parameters themselves are tunable, but the search adds compute and
%     is bounded by the same overfitting concerns AFML's CPCV is meant
%     to address.
%   - Limitation of both: walk-forward (or its equivalents) instead of
%     CPCV; neither computes DSR or PBO.
%   - This thesis fills the evaluation gap with full CPCV + DSR + PBO and
%     keeps TBL parameters fixed at AFML defaults to avoid the joint-
%     search overfitting risk.
%
% ── The fitting scheme matters more than model choice ─────────────
%
% Why this strand sits here, not in BTC ML: Chassot and Audrino (2026)
% test HAR vs. RF, lasso, GBT, FFNN on 1,445 US equities for realised
% volatility forecasting. Their target is not BTC, not direction, not
% classification; their claim is methodological.
%
% Topics:
%   - Across 1,445 US equities, a properly fitted linear HAR model
%     (Corsi 2009) outperforms RF, XGBoost, GBT, and FFNN for realised
%     volatility forecasting.
%   - Key finding: studies that report ML superiority over HAR used
%     suboptimal fitting schemes (infrequent re-estimation, short
%     training windows) that handicapped the baseline.
%   - When the fitting scheme is optimised (daily re-estimation,
%     2.5 to 4 year training window), the simpler model wins.
%   - Directly relevant to this thesis: the negative results (top DSR
%     0.2376, well below the 0.95 significance threshold) align with
%     this finding; properly evaluated baselines are hard to beat in
%     a weak-signal regime like BTC.
%   - Reinforces Q1 framing: methodology dominates model choice in
%     financial ML.
%
% Cite: \cite{slepaczuk_bieganowski_2024}, \cite{fu_et_al_2024},
% \cite{chassot_audrino_2026}, \cite{corsi_2009}.


\subsubsection{Gap statement}\label{subsubsec:afml_gap}

% Topics:
%   - Few crypto-ML studies adopt the full AFML pipeline; most cherry-
%     pick components.
%   - No study combines AFML and KAN for any asset.
%   - No study reports DSR and PBO for KAN-based predictions.
%   - This thesis is the first to apply the complete pipeline (CUSUM +
%     TBL + FFD + sample weights + CPCV + DSR + PBO) to BTC direction
%     prediction with KANs and five baselines.



% ==================================================================
% CHAPTER 3: METHODOLOGY
% ==================================================================
\newpage
\section{Methodology}\label{sec:methodology}

% Target: methodology + results + discussion together >= 70\% of textual
% part. Document the pipeline so a fellow graduate student can reproduce
% every number (Cochrane).
%
% Cochrane principle for this chapter: nothing before the main result
% that the reader does not need to understand the main result. Present
% sequentially; keep math to the minimum needed for results to be
% interpretable.
%
% T-R5: thread constructs (event, label, observation, fold, path)
% consistently throughout.
%
% Pipeline at a glance (Phase / Purpose / Output):
%   Phase I  (Pre-CPCV)    : Raw OHLCV to labelled events  -> (X, y, w, t1)
%   Phase II (CPCV pipeline): Per-fold training + prediction -> calibrated
%                              test-fold probabilities
%   Phase III (Post-CPCV)  : Evaluation + interpretability -> model
%                              comparison + symbolic formula


\subsection{Data}\label{subsec:data}

% Cochrane: do not write extensive descriptions of well-known datasets.
% Keep this section to about one page.
%
% Topics:
%   - Source: BTC-USD daily OHLCV from yfinance.
%   - Range: 2014-11-01 to 2026-05-01 (~4,200 daily bars).
%
% ── CUSUM start-date truncation ───────────────────────────────────
%
%   - The raw OHLCV series begins on 2014-11-01 to provide the 252-day
%     lookback the longest-warmup features require (Hurst at 252,
%     EMA 50/200 ratio at 200, SADF).
%   - The CUSUM event filter is then truncated to start on 2015-08-08,
%     the date of Ethereum's Frontier launch and the first day with
%     valid ETH/USD price data needed for the eth\_btc\_ratio feature.
%   - Events that fired before this date are dropped. The CUSUM
%     accumulators themselves are computed on the full raw series, so
%     events that survive truncation reflect the dynamic state of the
%     cumulative drift over the entire pre-event history.
%   - Empirically, the mean daily volatility over the truncated window
%     matches the full-series mean to four decimal places, confirming
%     the EWMA was fully converged by the truncation date.
%
% ── Why this buffer-and-truncate design ───────────────────────────
%
%   - First, every engineered feature must have completed its warmup
%     window before the first labelled event so that no fold sees a
%     partial-feature observation.
%   - Second, eth\_btc\_ratio cannot be computed before ETH started
%     trading; without truncation, early CPCV folds under N=8 produce
%     fully-NaN test partitions.
%   - Truncating CUSUM to start at the ETH availability date solves both
%     problems simultaneously.
%
% ── Validation pipeline ──────────────────────────────────────────
%
%   - Empty downloads raise.
%   - MultiIndex columns flattened.
%   - Duplicate dates raise.
%   - Calendar gaps <= 3 days forward-filled, gaps > 3 days raise.
%   - OHLCV consistency checks; NaN Close drops the row.
%
% ── Calendar ─────────────────────────────────────────────────────
%
%   - All rolling windows use the BTC trading calendar (7-day week,
%     30-day month, 90-day quarter, 180-day semester, 365-day year).
%   - BTC trades 24/7.
%
% ── External sources (overview, detailed in 3.4) ─────────────────
%
%   - Macro (yfinance, FRED).
%   - Crypto-macro (CoinMetrics with yfinance fallback).
%   - On-chain (CoinMetrics Community API).
%
% ── Anti-leakage ────────────────────────────────────────────────
%
%   - Externals aligned via merge\_asof(direction='backward').
%   - CoinMetrics shifted by 1 day (end-of-day reporting convention).
%
% ── Figure suggestion ────────────────────────────────────────────
%
%   - BTC-USD log price 2014 to 2026 with CPCV group boundaries
%     overlaid; vertical line marking the 2015-08-08 CUSUM truncation
%     date.


\subsection{Labelling}\label{subsec:labeling}

% T-R5: introduces the event / label / observation thread. Define the
% constructs in order of appearance:
%   - "event"       = CUSUM-triggered timestamp.
%   - "label"       = TBL output.
%   - "observation" = aligned final sample.


\subsubsection{Daily volatility}\label{subsubsec:vol}

% Topics:
%   - EWMA std of log returns, span = 50 (vs. AFML default 100 for
%     equities).
%   - Justification: BTC's faster regime transitions (halving cycles,
%     regulatory shocks).
%   - Two roles downstream: target width for TBL barriers; CUSUM
%     threshold calibration.
%
% Cite: \cite{lopez_de_prado_2018}, Snippet 3.1.


\subsubsection{CUSUM event filter}\label{subsubsec:cusum}

% Topics:
%   - Symmetric CUSUM accumulators s\_pos, s\_neg computed on the full
%     raw return series from 2014-11-01 onward.
%   - Threshold: h = 1.0 x mean(daily\_vol).
%   - Justification for 1.0x (tightened from 1.5x): empirical sweep
%     showed 0.5x too noisy, 3.0x too sparse; 1.0x yields workable
%     event count with balanced classes.
%   - Truncation: after CUSUM fires its candidate events on the full
%     series, the event index is truncated to start on
%     CUSUM\_START\_DATE = 2015-08-08 (Section 3.1). The accumulators
%     continue to reflect the dynamic state of the cumulative drift over
%     the entire pre-event history; only the event-firing window is
%     restricted to the data-availability frontier of eth\_btc\_ratio.
%   - Locked-run figures: 1,352 candidate events fire on the full series,
%     82 fall before the truncation date and are dropped, leaving 1,270
%     events that enter triple-barrier labelling.
%   - The post-truncation event series reduces 4,208 daily bars to 1,270
%     informative events; after triple-barrier labelling and rare-label
%     removal (Section 3.2.4) this becomes 1,168 binary-labelled
%     observations.
%
% Figure suggestion: CUSUM accumulators alongside BTC price with event
% markers; the 2015-08-08 truncation boundary marked.


\subsubsection{Triple barrier labelling}\label{subsubsec:tbl}

% Topics:
%   - pt\_sl = (1.5, 1.5) symmetric (no directional bias).
%   - num\_days = 10 (~ two trading weeks).
%   - min\_return = 0.02 (collapses small vertical-barrier returns into
%     class 0 for later removal).
%   - Output: DataFrame bins[ret, bin, t1] where t1 is the barrier-touch
%     timestamp (critical for downstream purging).
%   - Observed mean holding period 5.0 days (median 4 days); vertical
%     barrier hit on 19.7\% of events (horizontal barriers close
%     approximately 80\% before time barrier). Locked-run figures.
%
% Figure suggestion: triple-barrier visualization for two or three
% representative events.


\subsubsection{Rare label removal}\label{subsubsec:rare}

% Topics:
%   - min\_pct = 0.085 (raised from default 0.05 to aggressively remove
%     residual class-0 events).
%   - Combined with symmetric pt\_sl and min\_return = 0.02, class 0 is
%     eliminated -> binary labels {-1, +1}.
%   - Final aligned event count: 1,168 events after the 2015-08-08 CUSUM
%     truncation and the rare-class drop. Locked-run class balance:
%     \{+1: 664, -1: 504\} = 56.85\% Up / 43.15\% Down.
%   - Locked-run event-count chain: CUSUM fires 1,352 candidate events
%     on the full series; 82 fall before 2015-08-08 and are dropped by
%     the truncation, leaving 1,270 events that enter triple-barrier
%     labelling; rare-label removal drops 102 class-0 events (vertical-
%     barrier ties below the 0.02 minimum return threshold), giving the
%     final 1,168.
%   - The previous configuration (raw data from September 2014, no CUSUM
%     truncation) produced 1,245 events.
%
% Table suggestion: label distribution before and after rare-label
% removal. Self-contained caption.


\subsection{Sample Weights}\label{subsec:weights}

% Opening claim (T-R3): "Overlapping labels create redundancy that
% biases training; AFML's four-step weighting scheme corrects for it."
%
% Steps:
%   - Step 1. Concurrent label count (Snippet 4.1).
%       For each bar t, count active labels c\_t = |{i : t0\_i <= t <= t1\_i}|.
%   - Step 2. Average uniqueness (Snippet 4.2).
%       u\_bar\_i = mean(1 / c\_t) over [t0\_i, t1\_i]. Uniqueness near
%       1.0 -> highly informative; near 0 -> redundant.
%   - Step 3. Return-attribution weights (Snippet 4.10).
%       w\_i = |ret\_i| * u\_bar\_i, normalised to sum(w) = len(w)
%       (sklearn-compatible).
%   - Step 4. Time decay (Snippet 4.11).
%       Linear decay with time\_decay\_factor = 0.4 (oldest sample
%       weighted at 40\% of newest). Re-normalise.
%   - Step 5. Outlier capping.
%       Clip at weight\_cap\_quantile = 0.99.
%
% Integration:
%   - sklearn models: sample\_weight parameter in .fit().
%   - PyTorch models: per-sample multiplier in CrossEntropyLoss as
%     L = mean(w\_i * CE(logits\_i, y\_i)).
%
% Figure suggestion: sample weights over time, annotated with notable
% BTC events (March 2020 COVID, 2021 ATH, 2022 FTX collapse).
%
% Cite: \cite{lopez_de_prado_2018}, Chapter 4.


\subsection{Feature Engineering}\label{subsec:features}

% T-R4: one construct per paragraph; group, do not enumerate. Detailed
% per-feature table goes to Appendix E.
%
% Universe (73 features in four groups, all eligible for MDA selection;
% AR Logistic restricts itself to the 10 lag columns by name from the
% pre-MDA matrix):
%   - Technical (TA)               : 25  (OHLCV)
%   - Mathematical (AFML Part 4)   :  9  (returns / log-prices)
%   - External: macro              : 20  (yfinance, FRED)
%   - External: crypto-macro       :  1  (CoinMetrics + yfinance fallback)
%   - External: on-chain           :  8  (CoinMetrics)
%   - Lag (autoregressive)         : 10  (log returns)
%   - Total                        : 73


\subsubsection{Technical analysis features (25)}\label{subsubsec:ta}

% Group by function:
%   - Returns and volatility (8): log\_returns, realised vol
%     (annualised x sqrt(365)), Garman-Klass (1980), Yang-Zhang (2000),
%     ATR (EWMA, span=14, log-transformed), Bollinger Band width,
%     vol\_term\_7\_30 (7-day vs. 30-day realised vol ratio),
%     vol\_term\_30\_90 (30-day vs. 90-day realised vol ratio).
%   - Momentum and trend (9): RSI(14, Wilder), MACD/MACD-signal/
%     MACD-hist (12/26/9), roc\_14 (top BTC predictor in
%     mate\_confluence\_2024), Stoch \%K/\%D, Williams \%R, CCI(14).
%   - Volume (3): OBV (sign-preserving log-transformed), Chaikin
%     oscillator (3/10), MFI(14).
%   - Distribution shape (2): rolling skew/kurt (window=21).
%   - Trend ratios (3): EMA ratios 20/50, 50/200; VWMA ratio 20/50.
%
% Window convention: 21-day rolling for shape features; standard TA
% periods kept for named indicators (RSI=14, MACD=12/26/9). No
% optimisation of indicator periods (avoids another layer of overfitting
% risk).
%
% Cite: \cite{garman_klass_1980}, \cite{yang_zhang_2000},
% \cite{mate_confluence_2024}.


\subsubsection{Mathematical features (9, AFML Part 4)}\label{subsubsec:math}

% Topics:
%   - Information-theoretic (3): Shannon entropy (window=30), Lempel-Ziv
%     complexity (window=90), negentropy (window=30; difference between
%     the maximum-entropy Gaussian benchmark and the empirical Shannon
%     entropy of returns, so positive values flag departures from
%     Gaussianity).
%   - Random walk tests (2): Hurst (window=180, R/S at sub-windows
%     [10, 21, 42, 63]); variance ratio Lo-MacKinlay 1988 (window=90,
%     lag=7).
%   - Normality test (1): Jarque-Bera (window=90).
%   - Structural breaks (3): SADF (min sub-length=90, lags=1), SMT
%     polynomial-1, SMT exponential.
%   - Cached to cache/math\_features.parquet (O(n^2) for SADF and SMT).
%
% Cite: \cite{lopez_de_prado_2018} Part 4, \cite{lo_mackinlay_1988}.


\subsubsection{External features (29)}\label{subsubsec:external}

% Topics:
%   - Macro (20):
%       dxy\_roc\_30, us2y (FRED DGS2 with T10Y2Y fallback), us10y
%       (^TNX), yield\_curve\_2y10y, yield\_curve\_10y30y, vix, plus
%       30-day and 14-day return windows on each of seven commodity /
%       index series: sp500\_ret\_30 / sp500\_ret\_14, nasdaq\_ret\_30 /
%       nasdaq\_ret\_14, gold\_ret\_30 / gold\_ret\_14, silver\_ret\_30 /
%       silver\_ret\_14, copper\_ret\_30 / copper\_ret\_14, oil\_ret\_30 /
%       oil\_ret\_14, natgas\_ret\_30 / natgas\_ret\_14. The 14-day
%       variants were added to give the MDA selection step access to a
%       faster horizon on each commodity / index, since the previous
%       30-day-only configuration left no shorter macro horizon in the
%       MDA pool.
%
%   - Crypto-macro (1): eth\_btc\_ratio.
%       Computed as ETH\_close / BTC\_close aligned via merge\_asof.
%       ETH source priority: CoinMetrics Community API as the primary
%       source, trying three metrics in order (ReferenceRateUSD ->
%       PriceUSD -> derived CapMrktCurUSD / SplyCur); the first metric
%       returning more than 100 rows is used. yfinance ETH-USD is the
%       final fallback. This change replaces an earlier yfinance-only
%       implementation whose ETH-USD history began only in November 2017
%       and produced a 27\% NaN rate over the external dataframe with
%       entire-test-partition NaN in early CPCV folds under N=8. With
%       the CoinMetrics PriceUSD source, ETH coverage extends back to
%       2015-08-08 and the residual NaN rate falls to 7.7\%, all of
%       which sits in the November 2014 to August 2015 pre-ETH-trading
%       window and is fully truncated out by the CUSUM start-date filter
%       (Section 3.2.2). The earlier btc\_dominance column was removed
%       because the CoinGecko endpoint returns BTC market cap (not
%       bounded [0, 100] dominance) and the proxy fallback was a price-
%       correlated approximation that the methodology could not cleanly
%       defend.
%
%   - On-chain (8):
%       active\_addr\_roc\_14, tx\_count\_roc\_14, hashrate\_roc\_30,
%       mvrv (level), net\_exchange\_flow, fee\_per\_tx,
%       exchange\_supply\_pct, issuance\_ntv. CoinMetrics Community API,
%       shifted by 1 day.
%
%   - Anti-leakage: all external series merged onto BTC's calendar via
%     merge\_asof(direction='backward'). Cache invalidates on column-set
%     change (not just date range).


\subsubsection{Lag features (10)}\label{subsubsec:lags}

% Topics:
%   - AR\_LAGS = [1, 2, 3, 4, 5, 6, 7, 14, 21, 30], column prefix
%     log\_returns\_lag.
%   - The locked configuration adds four lags (4, 5, 6, 21) to the
%     previous six-lag set [1, 2, 3, 7, 14, 30]. The expansion fills the
%     gap between lag 3 and lag 7 with the missing trading-week pattern
%     and adds the three-week mark (lag 21) so AR Logistic has a full
%     near-term lag spectrum without the previous gap structure.
%   - Lag features compete with engineered features in MDA (advisor-
%     driven change to remove information asymmetry between AR Logistic
%     and the other models).
%   - AR Logistic still selects the 10 lag columns by name from
%     X\_tr\_full, regardless of MDA's choices.


\subsubsection{Alignment and NaN handling}\label{subsubsec:alignment}

% Three steps close out Phase 1: log transforms on a small set of
% unbounded features, NaN handling policy, and the final index
% alignment that produces (X, y, w, t1).
%
% ── Log transforms ────────────────────────────────────────────────
%
% Topics:
%   - atr: log(|x| + 1e-8).
%   - obv: sign(x) * log(|x| + 1).
%   - All other features are bounded or dimensionless.
%
% ── NaN handling across Phase 1 ───────────────────────────────────
%
% Topics:
%   - Phase 1 produces NaN values from rolling-window warmup and external
%     -data calendar gaps. Phase 1 does not drop or impute them.
%   - Alignment hard-asserts no entirely-NaN columns and fully populated
%     t1. Soft-warns on partial-NaN columns.
%   - Per-fold NaN resolution happens inside the CPCV loop
%     (Section 3.6.1).
%
% ── Alignment ────────────────────────────────────────────────────
%
% Topics:
%   - align\_for\_cv(features, bins, weights) -> (X, y, w, t1) via
%     index intersection.
%   - Hard assertions: non-empty intersection; no duplicate dates;
%     monotone index; identical lengths; no all-NaN columns.
%   - Validation: t1 fully populated; weights > 0; labels in {-1, 0, +1};
%     X/y/w/t1 share indices.
%   - Output: (1168 x 73) X, (1168) y, (1168) w, (1168) t1. Locked-run
%     class balance: \{+1: 664, -1: 504\} = 56.85\% Up / 43.15\% Down.
%
% Table suggestion: alignment summary (daily bars -> CUSUM events ->
% labelled events -> aligned size). Self-contained caption.


\subsection{Cross-Validation Framework}\label{subsec:cpcv}

% Opening claim (T-R3): "Standard k-fold CV is invalid for financial
% time series with overlapping labels; CPCV prevents the resulting
% leakage." Invest 1.5 to 2 pages. The committee must understand purging
% and embargo to evaluate the results.


\subsubsection{Why standard CV fails}\label{subsubsec:cv_fail}

% Topics:
%   - TBL labels span [t0, t1] (multi-day intervals).
%   - Overlapping spans across folds -> training labels carry information
%     about test-period prices.
%   - Inflates metrics: model appears to predict the future but is
%     recognising patterns from overlapping training labels.


\subsubsection{CPCV configuration}\label{subsubsec:cpcv_config}

% Topics:
%   - N\_GROUPS = 8, K\_TEST\_GROUPS = 2, EMBARGO\_PCT = 0.01.
%   - Groups 0 to 6 of size floor(T/N), group 7 absorbs remainder.
%   - Total splits: C(8, 2) = 28; backtest paths:
%     C(N-1, k-1) = C(7, 1) = 7.
%   - Each group appears in 7 test sets. Per-group sample size
%     approximately 146 events at the locked configuration.
%
% Justification for N=8, k=2:
%   - Yields approximately 146 events per group while keeping the
%     training fold at approximately 858 events on average (73.4\% of
%     the 1,168 aligned events) after purging and embargo.
%   - 28 splits and 7 paths give denser combinatorial diversity for PBO
%     than the earlier N=6 configuration (which produced 15 splits and
%     5 paths) without dropping per-group sample size below the rough
%     lower bound for daily-bar AFML pipelines.
%   - The choice trades a smaller test fold per split against a larger
%     Sharpe-matrix cross-section for PBO and DSR.
%
% Locked-run audit (informational): across all 28 splits the average
% train size is 858 events, average test size is 292, average purged-
% per-split is 0.8 observations, and average embargoed-per-split is
% 17.5 observations. Per-group sample size is 144 for groups 0 to 6
% and 151 for the remainder group 7. Group date ranges span 2015-08-08
% (G0) to 2026-04-19 (G7). Zero leakage detected across all 28 splits.
%
% Table suggestion: group boundaries (group ID, positional index range,
% date range, count). Self-contained caption.


\subsubsection{Purging, embargo, and leakage verification}\label{subsubsec:purging_embargo_leakage}

% Three mechanisms close the leakage channels that overlapping TBL
% labels create across CPCV folds: purging removes overlapping training
% labels, embargo absorbs serial correlation past the test boundary,
% and an empirical leakage audit verifies both work.
%
% ── Purging ───────────────────────────────────────────────────────
%
% Three sufficient overlap conditions for training observation i against
% test [t\_test\_start, t\_test\_end]:
%   1. t\_test\_start <= t0\_i <= t\_test\_end
%      (observation falls in test window).
%   2. t\_test\_start <= t1\_i <= t\_test\_end
%      (label resolves in test window).
%   3. t0\_i <= t\_test\_start AND t\_test\_end <= t1\_i
%      (label spans the entire test window).
%
% Any training observation satisfying at least one condition is removed
% for that split.
%
% Cite: \cite{lopez_de_prado_2018}, Snippet 7.1.
%
% ── Embargo ──────────────────────────────────────────────────────
%
% Topics:
%   - int(EMBARGO\_PCT * T) = 11 observations removed immediately after
%     each test group (1.0\% of T=1,168, rounded down).
%   - Applied only after the test group, not before: training labels
%     resolving before the test starts contain no future test
%     information.
%   - Prevents serial-correlation leakage (autocorrelation in volatility,
%     momentum spillover).
%
% Cite: \cite{lopez_de_prado_2018}, Section 7.4.2.
%
% ── Leakage verification ─────────────────────────────────────────
%
% T-R15: write for reviewers.
%
% Topics:
%   - Empirical audit across all 28 splits: zero training observations
%     whose t1 falls within the assigned test group's date range.
%
% Table suggestion: leakage audit (28 rows, columns: split, test groups,
% train count, test count, leaks, status). Self-contained caption.


\subsubsection{Path matrix}\label{subsubsec:path_matrix}

% Topics:
%   - 7 backtest paths, each covering all 8 groups exactly once.
%   - For group g, path p uses the p-th split that includes g in its
%     test set.
%   - Enables distribution-level Sharpe analysis instead of single-point
%     estimates.
%
% Figure suggestion: CPCV split visualization for representative splits
% (0, 13, 27) with train, test, purged, and embargo regions.
%
% Cite: \cite{lopez_de_prado_2018}, Section 12.4.1.


\subsection{Per-Fold Preprocessing}\label{subsec:preprocessing}

% Opening claim (T-R3): "Three transformations happen inside the CPCV
% loop, fitted on training data only, to prevent test-fold statistics
% from leaking into the training pipeline."


\subsubsection{Fractional differentiation}\label{subsubsec:ffd}

% Topics:
%   - ADF test across all 73 features at alpha = 0.05 identifies ATR as
%     the only non-stationary feature. FFD applied to ATR only.
%
% Procedure:
%   - Weights: omega\_0 = 1, omega\_k = -omega\_{k-1} * (d - k + 1) / k;
%     truncate when |omega\_k| < 1e-4 or k >= 200.
%   - d* sweep: d in [0, 1], step 0.05; minimum d where ADF p-value
%     < 0.05.
%   - If no d achieves stationarity, default to d* = 1.0.
%
% Application scope:
%   - FFD applied to the full ATR series using the training-derived d*,
%     so test observations have valid lookback history. Strictly
%     backward-looking convolution.
%
% Per-fold NaN policy (asymmetric by column type):
%   - Non-FFD columns: ffill().bfill() independently within X\_train and
%     within X\_test (no boundary crossing).
%   - FFD columns: drop NaN rows (forward-filling FFD output would
%     inject stale lookback values). Logs "FFD: dropped X train, Y test
%     NaN rows from FFD lookback." Typically 2 to 10 rows.
%
% Cite: \cite{lopez_de_prado_2018}, Ch. 5; \cite{slepaczuk_bieganowski_2024}
% for walk-forward FFD validation.


\subsubsection{Feature scaling}\label{subsubsec:scaling}

% Topics:
%   - RobustScaler (median + IQR), fitted on training fold only, applied
%     to all 73 features.
%   - Choice over StandardScaler: BTC features show fat tails and
%     extreme values; median/IQR resists outliers, mean/std is heavily
%     influenced by them.


\subsubsection{Multi-model MDA feature selection}\label{subsubsec:mda}

% Topics:
%
%   Multi-model design (novel relative to single-model MDA):
%   - MDA computed independently with Random Forest (500 trees, balanced
%     class weights, captures nonlinear interactions).
%   - MDA computed independently with Logistic Regression (balanced,
%     captures linear effects).
%   - Final MDA = mean(MDA\_RF, MDA\_LR) per feature.
%   - Rationale: prevents bias toward any single model architecture; SFI
%     in weak-signal regimes returns near-uniform scores; RF-only
%     inflates tree-friendly features.
%
%   Inner CV: purged 3-fold on the training set (same t1-based overlap
%   conditions as outer CPCV).
%
%   Selection rule: keep features with averaged MDA > 0; cap at
%   MDA\_TOP\_K\_FRAC = 0.30 (approximately 22 of 73); minimum floor of
%   5 features.
%
%   TOP\_K\_FRAC tightening rationale (advisor-reviewed):
%   - The cap was tightened from an earlier 0.40 to 0.30 in the locked
%     configuration after a high-PBO run with the looser setting.
%   - Across the earlier-run CPCV folds, only approximately 6 features
%     cleared 50\% selection frequency in the stability bar chart,
%     indicating that the long tail of the MDA-ranked feature set was
%     contributing variance rather than signal.
%   - Tightening to 0.30 forces approximately 22 features through the
%     bottleneck and aligns the selection cap with the empirical
%     stability finding.
%   - Trade-off: one or two folds may select fewer features than they
%     would have at 0.40, but those folds were also the ones contributing
%     the most rank variance to PBO, so the tightening attacks the right
%     problem.
%   - An even tighter 0.25 cap (approximately 18 features) was tested in
%     an interim configuration; 0.30 was retained as the final value
%     because it preserves enough feature diversity in the inner Optuna
%     step without re-introducing the long tail.
%
%   AR Logistic exception: bypasses MDA entirely; receives pre-MDA
%   matrix and selects 10 lag columns by name.
%
%   Typical result: approximately 22 features selected per fold from 73
%   candidates (down from approximately 26 to 30 under the previous 0.40
%   cap; up from approximately 18 under the interim 0.25 cap).
%
% Cite: \cite{lopez_de_prado_2018}, Chapter 8 (Section 8.4).


\subsection{Models}\label{subsec:models}

% T-R4: one construct per paragraph; six models, four families. Six
% per-family subsections describe what is unique about each model,
% followed by a Shared neural design subsection and a summary table as
% a closer. The framework-level Optuna details (TPE sampler, inner CV,
% trial budget, MedianPruner, per-split tuned-params application) live
% in the separate Hyperparameter Tuning subsection (Section 3.8); this
% section keeps the per-model "Tuned per split:" bullets that list the
% search spaces alongside each model's fixed-architecture choices, but
% does not restate the framework mechanics.
%
% Note for the writing pass: the per-model "Tuned per split:" bullets
% list the search spaces; the per-model fixed-architecture choices are
% above them. Anything about how Optuna explores those search spaces
% (TPE, pruner, trial budget, parameter application) belongs in 3.8,
% not here.
%
% Shared elements that live OUTSIDE this chapter: sample weights are
% in Section 3.3; class balancing is in each per-model subsection
% below (via class\_weight='balanced' or scale\_pos\_weight); per-fold
% calibration is in Section 3.9 (renumbered after the new 3.8 split:
% Models 3.7, Hyperparameter Tuning 3.8, Calibration 3.9, Pipeline
% Orchestration 3.10, Evaluation Framework 3.11, Symbolic Extraction
% 3.12, Leakage Prevention Summary 3.13; the previous 3.13 Diagnostics
% chapter is removed and its calibration-audit content is noted in
% Section 3.9 as the diagnostic that exposed the vector-scaling issue).


\subsubsection{AR Logistic (econometric baseline)}\label{subsubsec:ar}

% Topics:
%   - Tests pure price momentum vs. 73 engineered features.
%   - Lags [1, 2, 3, 4, 5, 6, 7, 14, 21, 30] of log returns.
%   - Architecture: sklearn LogisticRegression with C=1.0, L2,
%     class\_weight='balanced', solver='lbfgs', max\_iter=1000.
%   - 5 seeds, no per-split tuning (deterministic baseline). The locked configuration uses 5 seeds uniformly across all six models (up from 3 in the previous run).
%   - Consumes the 10 lag columns by name from X\_tr\_full, independent
%     of MDA.


\subsubsection{Logistic Regression (linear ML baseline)}\label{subsubsec:lr}

% Topics:
%   - On MDA-selected features.
%   - class\_weight='balanced' + AFML sample\_weight (dual weighting).
%   - Tuned per split: C (log-uniform [1e-4, 1e2]), penalty in {l1, l2}.
%   - Solver auto-selected: liblinear for L1, lbfgs for L2.
%   - 5 seeds, 30 trials per split.


\subsubsection{Random Forest}\label{subsubsec:rf}

% Topics:
%   - 500-tree ensemble, class\_weight='balanced\_subsample', n\_jobs=-1.
%   - Tuned per split: n\_estimators in [100, 250] step 50 (capped from
%     an earlier 300 ceiling; trees in a noisy regime do not benefit from
%     more than 250); max\_depth in [2, 6] (tightened from earlier
%     [3, 15]; depth 6 has 64 leaves which is plenty for approximately
%     851-sample training folds, and shallower forests vote in tighter
%     agreement, reducing the disagreement that surfaces as path-Sharpe
%     variance); min\_samples\_leaf in [15, 40] (raised from earlier
%     [1, 30]; a floor of 15 forces each leaf to represent at least
%     1.8\% of the training fold, preventing leaves that fit just a
%     handful of high-volatility events); max\_features in {sqrt, log2}.
%   - 5 seeds, 30 trials per split.
%
% Cite: \cite{breiman_2001}.


\subsubsection{XGBoost}\label{subsubsec:xgb}

% Topics:
%   - 500-tree gradient-boosted ensemble with early stopping at 20
%     rounds.
%   - Objective binary:logistic; scale\_pos\_weight from class balance.
%   - Tuned per split: max\_depth in [1, 3] (tightened from earlier
%     [2, 6]; XGBoost's sequential boosting compounds depth nonlinearly
%     across rounds, so depth 3 across 50 boosting rounds already
%     produces substantial nonlinear capacity, and depth 6 in this regime
%     memorises residuals); learning\_rate log-uniform [0.01, 0.3] (floor
%     at 0.01; below this, training takes forever and effectively
%     underfits); min\_child\_weight in [5, 30] (floor raised from 1 to
%     align with RF's leaf-size discipline; with approximately 858 train
%     samples a min\_child\_weight=1 permits trees to split off single-
%     event leaves); subsample and colsample\_bytree in [0.6, 1.0];
%     gamma log-uniform [1e-8, 1.0]; reg\_alpha, reg\_lambda log-uniform
%     [1e-8, 10.0].
%
%   - Calibration set dual role: calibration set acts as eval set for
%     early stopping AND as Platt-fit data. Acknowledged as a mild
%     dependency: only ensemble size affected, no individual tree
%     decisions.
%
%   - 5 seeds, 30 trials per split.
%
% Cite: \cite{chen_guestrin_2016}.


\subsubsection{LSTM}\label{subsubsec:lstm}

% Methodology assumes the reader has read 2.1; do not re-explain the
% LSTM architecture from first principles.
%
% Topics:
%   - Architecture: single-layer nn.LSTM (num\_layers=1 hardcoded) ->
%     last hidden state from the final layer -> LayerNorm -> dropout ->
%     linear classifier. Hidden size, dropout, and learning rate are
%     tuned per split; num\_layers is no longer searched.
%   - Sliding window: LSTM\_WINDOW = 14 (deliberately close to TBL
%     num\_days = 10; longer windows attenuate gradient signal and
%     inflate parameter-to-sample ratio). Reduces effective training
%     count from N to N - 13 sequences. last\_valid\_indices stored for
%     re-alignment.
%   - Last-hidden-state pooling: earlier learned-attention pooling was
%     removed; with window=14 and approximately 858-sample folds,
%     additional attention parameters did not improve performance.
%   - Tanh input normalisation: z = tanh((x - mu) / sigma), mean and
%     std fitted on training data only.
%   - Training stack: AdamW (lr tuned, weight\_decay=1e-4),
%     CrossEntropyLoss with class weights and AFML sample weights, label
%     smoothing 0.1, gradient clipping (max norm 1.0), cosine annealing
%     warm restarts (T\_0=25, T\_mult=2), batch size 64, max 100 epochs,
%     early stopping patience 15, best-state restoration.
%   - Tuning consistency: LSTMClassifier.\_\_init\_\_ reads module-level
%     constants at call time (not as default args), so tuning overrides
%     actually reach the model. Tuning runs at epochs=50, patience=7;
%     production refits at epochs=100, patience=15. This is the only
%     axis where tuning and production diverge; documented as a
%     deliberate compute-vs-fidelity trade-off.
%   - Tuned per split: hidden\_size in [16, 32] step 16; num\_layers
%     fixed at 1 (no longer searched; tightened from earlier [1, 2] then
%     [1, 3]; two- and three-layer LSTMs on 1,168 events are deep-overfit
%     territory and the additional layer added variance to path-Sharpes
%     without improving accuracy. Hardcoding to 1 frees Optuna trials for
%     finer exploration of dropout and learning\_rate); dropout in
%     [0.1, 0.5] (floor raised from 0.0); lr log-uniform [1e-4, 5e-2].
%   - 5 seeds, 30 trials per split (was 2 seeds in the previous run; the locked configuration now uses 5 seeds uniformly across neural and classical models).
%
% Cite: \cite{hochreiter_schmidhuber_1997}, \cite{loshchilov_adamw_2019}.


\subsubsection{KAN}\label{subsubsec:kan_model}

% Methodology assumes the reader has read 2.1; this section covers
% implementation specifics only (efficient-kan vs PyKAN, this thesis's
% architecture, hyperparameters).
%
% Topics:
%   - Library and architecture: efficient\_kan.KAN([n\_features, width1,
%     2], grid\_size=grid, spline\_order=3, grid\_range=[-1, 1]).
%     Single hidden layer by construction: the second hidden layer is
%     permanently disabled (width2=0 hardcoded in tuning.py). Width is
%     tuned per split (width1 in [2, 6]) and grid size (in {3, 5}).
%
%   - Why single-hidden-layer architecture:
%       The CPCV-evaluated KAN matches the architecture used in the
%       Phase 3 symbolic extraction. The benchmark numbers and the
%       extracted symbolic formula therefore describe the same model
%       rather than two unrelated KAN topologies. A two-hidden-layer
%       KAN would let symbolic extraction nest trigonometric primitives
%       in trigonometric primitives, producing fourth-order compositions
%       that lose interpretability and would force the symbolic chapter
%       to caveat that the formula approximates a different model than
%       the one benchmarked.
%
%   - Tanh input normalisation matching grid range.
%   - Training stack: AdamW (lr and weight\_decay tuned),
%     CrossEntropyLoss with class weights and AFML sample weights,
%     label smoothing 0.1, gradient clipping (max norm 1.0), cosine
%     annealing warm restarts (T\_0=30, T\_mult=2), early stopping
%     patience 20, best-state restoration. Max 200 epochs.
%   - Single-grid training (no coarse-to-fine): with approximately 858
%     training samples per CPCV fold, grid refinement adds parameters
%     faster than the data can support.
%   - No SWA, no entropy regularisation: SWA conflicted with early
%     stopping; entropy regularisation was redundant with
%     label\_smoothing=0.1. Removed for coherence.
%
%   - Dual-library strategy: efficient-kan for all 28 CPCV splits (fast,
%     stable, standard PyTorch). PyKAN re-trained independently for
%     symbolic extraction (Section 3.11), where only PyKAN exposes
%     prune(), suggest\_symbolic(), fix\_symbolic(), symbolic\_formula().
%     Both share the B-spline basis, tanh normalisation, and now the
%     single-hidden-layer topology.
%
%   - Tuned per split: width1 in [2, 6] (tightened from earlier [3, 12]
%     then [3, 16]; the cap is set at 6 to keep the symbolic formula
%     extracted in Phase 3 humanly readable, since each surviving width1
%     unit becomes one additive term plus interactions in the closed-
%     form expression); width2 fixed at 0 (no longer searched);
%     lr log-uniform [5e-4, 5e-2]; weight\_decay log-uniform
%     [1e-5, 5e-3]; grid in {3, 5} (dropped grid=8 to prevent
%     memorisation).
%   - 5 seeds, 30 trials per split (was 2 seeds in the previous run; uniform 5-seed configuration across all six models in the locked v6 run).
%
% Cite: \cite{liu_kan_2024}.


\subsubsection{Shared neural design}\label{subsubsec:shared_neural}

% Topics:
%   - LSTM and KAN both use dual weighting in CrossEntropyLoss: class
%     weights (inversely proportional to class frequency) AND AFML
%     sample weights as per-sample multipliers.
%   - Uniform 5-seed configuration (v6 locked run): all six models
%     use 5 seeds. The previous run used 3 seeds for classical models
%     and 2 for neural models (an asymmetric configuration justified
%     by the 5x to 10x cost premium of neural training); the current
%     locked run pays the additional compute cost to remove the
%     asymmetry. Split-level metrics are averaged across 5 seeds for
%     every (model, split) cell, and the evaluation code handles the
%     uniform structure straightforwardly. The previous run's seed
%     asymmetry was the largest single configuration difference
%     between v5 and v6.
%   - All models share the BaseModel interface: fit, predict\_proba,
%     predict\_logits, get\_name. Identical training/evaluation
%     conditions across architectures.


\subsubsection{Summary table}\label{subsubsec:model_summary}

% Self-contained caption (Cochrane): "Summary of the six models
% evaluated under CPCV (N=8, k=2, 28 splits, 7 backtest paths). All
% models receive AFML sample weights and balanced class weights.
% Hyperparameters listed under 'Tuned' are optimised per fold via Optuna
% TPE (Section 3.8); 'Fixed' parameters are held constant across all
% folds. Role describes what hypothesis each model tests relative to
% the research questions."
%
% Columns: Model | Family | Architecture | Fixed params | Tuned params |
%          Seeds | Trials | Role
%
% Rows (use locked-run tuning ranges, NOT the old [3,12] / [1,3] /
% [3,15]):
%   AR Logistic           | Econometric | LR on lags [1,2,3,4,5,6,7,14,21,30]   | C=1.0, L2, max\_iter=1000   | (none)                                                                                              | 5 | 0  | Pure momentum (Q1, Q2)
%   Logistic Regression   | Linear ML   | LR on selected features                | max\_iter=1000              | C, penalty                                                                                          | 5 | 30 | Linear baseline (Q3)
%   Random Forest         | Ensemble    | balanced\_subsample                    | n\_jobs=-1                  | n\_estimators <= 250, max\_depth in [2, 6], min\_samples\_leaf in [15, 40], max\_features            | 5 | 30 | Nonlinear ensemble (Q3)
%   XGBoost               | Ensemble    | 500 trees + early stop@20              | binary:logistic             | max\_depth in [1, 3], lr, min\_child\_weight in [5, 30], subsample, etc. (8)                        | 5 | 30 | Gradient boosting (Q3)
%   LSTM                  | Neural      | window=14, last-hidden, num\_layers=1 | T\_0=25, batch=64           | hidden in {16, 32}, dropout in [0.1, 0.5], lr                                                       | 5 | 30 | Temporal dependencies (Q3)
%   KAN                   | Neural      | [n, width1, 2], single hidden, k=3     | width2=0, T\_0=30, ls=0.1   | width1 in [2, 6], grid in {3, 5}, lr, weight\_decay                                                 | 5 | 30 | Interpretable architecture (Q3, Q4)


\subsection{Hyperparameter Tuning}\label{subsec:tuning}

% Opening claim (T-R3): "Hyperparameter selection happens entirely
% inside each CPCV training fold; the framework here describes the
% algorithm that explores those search spaces, not the per-model
% ranges (which live in Section 3.7)."
%
% Topics:
%   - Architecture: nested per-split Optuna study inside each CPCV
%     training fold. AFML Ch. 7 compliant: outer CPCV provides unbiased
%     test folds; inner tuning operates entirely within the training
%     fold. Test fold never seen during tuning.
%   - Inner CV: purged 3-fold (N\_INNER\_FOLDS=3), 10-observation
%     embargo around inner-fold boundaries (matches TBL num\_days=10).
%   - Sampler and pruner: TPE sampler (seed=42) drives exploration;
%     MedianPruner with n\_startup\_trials=5 for classical models and
%     3 for neural models, n\_warmup\_steps=1. The pruner kills trials
%     whose intermediate log loss falls below the median of completed
%     trials, freeing compute for promising configurations.
%
%   - Trial budget:
%     n\_trials=30 per tuned model per split, applied uniformly across
%     LR, RF, XGBoost, LSTM, and KAN. The notebook passes a single
%     n\_trials=30 parameter so every tuned model competes on the same
%     budget; AR Logistic is not tuned. The wider module defaults
%     (N\_TRIALS\_CLASSICAL=60, N\_TRIALS\_NEURAL=40) are reserved for
%     sensitivity experiments outside the headline run.
%
%   - Per-split tuned-params application mechanism:
%     \_apply\_tuned\_params writes the best params to module-level
%     constants (e.g., kan\_model.KAN\_HIDDEN) before the model
%     training loop for that split. \_reset\_module\_defaults snapshots
%     pristine values on first invocation and restores them on
%     subsequent runs to prevent contamination across calls.
%
%   - DSR/PBO validity: n\_trials in DSR counts the number of compared
%     models (6), NOT the Optuna trials per split. Optuna trials happen
%     inside the training fold and do not affect test-fold Sharpe
%     estimates.
%
% Cite: \cite{akiba_optuna_2019}, \cite{lopez_de_prado_2018} Ch. 7.


\subsection{Calibration}\label{subsec:calibration}

% Opening claim (T-R3): "Bet sizing depends on calibrated probabilities;
% miscalibrated probabilities produce systematically wrong position
% sizes."
%
% Topics:
%   - Calibration set: 80/20 chronological split of the training fold;
%     calibrator fitted on the held-out 20\% (approximately 170 obs at
%     N=8, with model-training set approximately 680 obs after the
%     split). Never touches test data.
%
% Two methods, auto-selected by model type:
%   - Platt scaling (Platt 1999): AR Logistic, LR, RF, XGBoost. Input:
%     1D log-odds. Mechanism: LogisticRegression(C=1e10) mapping logits
%     -> calibrated proba.
%   - Vector scaling (Guo et al. 2017, Section 4.2): LSTM, KAN. Input:
%     2D logits. Mechanism: fit T and per-class bias b minimising NLL of
%     softmax((logits + b) / T) via L-BFGS-B with T in [0.05, 20] and
%     b\_c in [-5, 5].
%
% Why vector scaling rather than temperature scaling (T-R15):
%   - A pre-final calibration audit revealed that LSTM and KAN
%     systematically under-predicted P(y=1) by 10 to 23 percentage
%     points, while the empirical base rate of class 1 was approximately
%     0.55.
%   - Pure temperature scaling preserves the argmax of raw logits by
%     construction, so a single T cannot shift "lean class 0" to "lean
%     class 1".
%   - The bias propagated through bet sizing as systematic short bets in
%     upward-drifting regimes, contributing to a negative path Sharpe in
%     the early KAN equity curves.
%   - Vector scaling adds a per-class bias b that lifts the directional
%     constraint (Guo et al. 2017 recommends this as the natural
%     extension when temperature scaling alone is insufficient).
%   - The substitution was made before the final evaluation pass: a
%     correction of methodological inadequacy, not test-set-informed
%     model selection.
%
% Calibration set dual role disclosure:
%   - The 20\% subset serves as both early-stopping monitor (XGBoost)
%     and Platt/vector input.
%   - Each calibration method fits at most three parameters; early
%     stopping affects only ensemble size, not individual tree decisions.
%   - Splitting the small cal set further would degrade both purposes.
%
% Calibration audit (origin of the vector-scaling choice): a binned
% predicted-vs-empirical diagnostic, computed on pooled CPCV test-fold
% predictions (bins >= 10 samples), is what surfaced the systematic
% LSTM/KAN under-prediction of P(y=1) described above and motivated
% the substitution of vector scaling for temperature scaling. The
% same diagnostic is now part of the standard post-run reporting.
%
% Cite: \cite{platt_1999}, \cite{guo_temperature_2017}.


\subsection{Pipeline Orchestration}\label{subsec:pipeline}

% Topics:
%   - Single entry point: run\_cpcv\_pipeline(X, y, w, t1, bins\_ret,
%     ..., tune=False, n\_trials=None).
%   - Module-default reset on every invocation prevents tuned-value
%     contamination across runs.
%
% Per-split execution (28 splits total):
%   1. Extract fold data via positional indices.
%   2. preprocess\_fold() (FFD -> scale -> MDA select).
%   3. Re-align y, w, t1 after FFD may drop rows.
%   4. Keep X\_tr\_full for AR Logistic alongside X\_tr\_sel for the
%      others.
%   5. Chronological 80/20 train/cal split.
%   6. If tune=True: nested Optuna on X\_tr\_sel, apply best params to
%      module constants.
%   7. For each (model, seed): create model, route correct X, fit,
%      calibrate, predict on test.
%   8. Store keyed by (model\_name, split\_idx, seed).
%
% Models run separately in the notebook (Section 4). Results merged via
% dict union. Failed fits caught, logged, skipped.


\subsection{Evaluation Framework}\label{subsec:evaluation}

% Module constants:
%   - TRANSACTION\_COST = 0.001
%   - MIN\_BET\_SIZE = 0.05
%   - MAX\_BET\_SIZE = 0.75
%   - BET\_DISCRETIZATION = [0.0, 0.25, 0.50, 0.75]   (1.0 dropped,
%     consistent with the 0.75 cap)
%   - ANNUALIZATION\_FACTOR = 365
%   - RISK\_FREE\_RATE = 0.0


\subsubsection{Trading simulation}\label{subsubsec:trading_simulation}

% This subsection covers the full bet-sizing-to-path-stitching loop:
% how calibrated probabilities become directional bets, how bets
% become strategy returns, and how the 28 per-split prediction streams
% get assembled into 7 backtest paths.
%
% ── Bet sizing ────────────────────────────────────────────────────
%
% Topics:
%   - Direction: sign = +1 if P(up) > P(down) else -1.
%   - Confidence: p = max(P(0), P(1)) in [0.5, 1.0].
%   - Z-score: z = (p - 0.5) / sqrt(p * (1 - p) + 1e-10).
%   - Raw bet: 2 * Phi(z) - 1.
%   - Cap: clip(raw\_bet, -0.75, 0.75).
%   - Threshold: |bet| < 0.05 -> 0.
%   - Discretise: nearest of {0, 0.25, 0.5, 0.75}.
%   - Apply sign.
%
% Implicit abstention mechanism: predictions with p ~ 0.50 produce bet
% ~ 0; aligns with the Conservative Predictions framework
% (Section 2.2.4).
%
% ── Strategy returns ─────────────────────────────────────────────
%
% Topics:
%   - gross = bet * label\_return; turnover = |Delta bet|;
%     tx\_cost = 0.1\% * turnover; net = gross - tx\_cost.
%   - 0.1\% per-trade cost: conservative for major BTC exchanges
%     (maker 0.02\% to 0.10\%, taker 0.04\% to 0.10\%, plus slippage).
%   - Annualisation at 365 (BTC trades 24/7); risk-free rate 0
%     (no universally accepted crypto risk-free rate).
%
% ── Path stitching ───────────────────────────────────────────────
%
% Topics:
%   - 7 backtest paths assembled from the 28 splits.
%   - Collects (group\_id, split\_id) from path\_map[path\_id].
%   - For each pair: retrieves the split's stored predictions AND
%     filters to events whose positional index falls within
%     group\_bounds[group\_id]. Critical: without the filter, events
%     from co-tested groups get duplicated.
%   - Multi-seed: averages calibrated probabilities across seeds before
%     bet sizing (variance reduction by approximately 1/sqrt(n\_seeds)).
%   - Concatenate, sort, assert no duplicate timestamps.
%
% Bug-fix disclosure (T-R15):
%   - Earlier implementation pulled each split's full test set, double-
%     or quintuple-counting events.
%   - Surfaced as a duplication pattern by direct timestamp inspection.
%   - Fix is the group filter described above. All path-level metrics in
%     this thesis use the corrected stitching; the bug is preserved here
%     as a transparency disclosure.
%
% Cite: \cite{lopez_de_prado_2018} Ch. 10 and Section 12.4.1;
% \cite{nabar_shroff_2023}.


\subsubsection{Path performance metrics}\label{subsubsec:path_metrics}

% Metrics computed per path:
%   - Annualized Sharpe: (mean / std) * sqrt(365) (daily-equivalent
%     scale).
%   - Cumulative return: prod(1 + r\_t) - 1.
%   - Annualized return: (1 + cum\_ret)^(1 / years\_elapsed) - 1
%     (calendar-time CAGR).
%   - Max drawdown: min((equity - running\_max) / running\_max).
%   - Time under water: longest consecutive run below previous peak
%     (days).
%   - Win rate: fraction of traded observations with positive return.
%   - Profit factor: sum(positive) / |sum(negative)|; NaN if no trades,
%     inf if no losses.
%   - n\_trades: count where bet != 0.
%   - n\_returns: length of strategy returns including zero-bet rows.
%   - Skew, kurt: distribution shape of strategy returns.
%
% Sharpe annualisation disclosure: strategy returns are sampled at CUSUM
% events (approximately 75 per year), not daily bars (365 per year). The
% sqrt(365) choice yields the daily-equivalent scale; bet-frequency
% annualisation sqrt(75) ~ 8.7 would give roughly half the reported
% Sharpe. Internally consistent (DSR de-annualises by the same
% sqrt(365), so DSR verdicts and rankings are convention-invariant), but
% the absolute Sharpes sit on the daily-equivalent scale. Disclosed
% explicitly.


\subsubsection{Statistical correction layer: DSR, PBO, and DeLong}\label{subsubsec:correction_layer}

% Three corrections are applied jointly: DSR for multiple-testing on the
% Sharpe, PBO for IS/OOS overfitting on the selection process, and
% DeLong for AUC differences across model pairs.
%
% ── Deflated Sharpe Ratio ────────────────────────────────────────
%
% E[max SR] = sqrt(2 * ln(n)) * (1 - gamma / (2 * ln(n))) +
%             gamma / (2 * sqrt(2 * ln(n)))
% SR\_std    = sqrt((1 - skew*SR + (kurt + 2)/4 * SR^2) / (n\_obs - 1))
% DSR       = Phi((SR\_observed - E[max SR]) / SR\_std)
%
% Topics:
%   - n = n\_trials = 6 (number of compared models). NOT Optuna trials.
%   - Kurtosis convention: Mertens (2002) formula assumes raw kurtosis
%     (3 for normal); scipy.stats.kurtosis returns excess (0 for normal).
%     Implementation converts: (gamma\_4 - 1)/4 = (excess + 2)/4. For
%     normal returns reduces to 1 + SR^2/2 per Lo (2002).
%   - n\_obs source: compute\_model\_summary passes avg\_n\_returns
%     (mean across paths of len(strategy\_returns)), matching the n that
%     estimated the Sharpe.
%   - NaN-safety clamp on the variance term before sqrt.
%   - DSR > 0.95 -> significant after multiple-testing correction.
%
% ── Probability of Backtest Overfitting (CSCV) ───────────────────
%
% Topics:
%   - Path Sharpe matrix shape (6 models x 7 paths).
%   - All C(7, 3) = 35 IS/OOS partitions of the 7 paths (IS=3 paths,
%     OOS=4 paths; with 7 paths the split is asymmetric by construction).
%   - For each partition: identify IS-best, check if it underperforms
%     OOS median.
%   - PBO = fraction of partitions where IS-best underperforms.
%   - PBO < 0.3 -> robust selection; PBO > 0.5 -> anti-predictive (IS
%     winner is OOS loser).
%   - The N=8 configuration delivers 35 IS/OOS partitions vs. 10 under
%     N=6, increasing PBO resolution from 0.10 to roughly 0.029 per
%     partition.
%
% ── DeLong pairwise AUC ──────────────────────────────────────────
%
% Topics:
%   - Per (model, split): average predicted probas across the 5 seeds
%     (uniform across all six models in the locked v6 configuration;
%     was 3 in the previous run).
%     Matches what stitch\_paths does for path-level metrics.
%   - Pool seed-averaged predictions across all 28 splits per model
%     (valid because CPCV test sets are non-overlapping).
%   - Z-statistic: z = (AUC\_a - AUC\_b) / sqrt(Var(AUC\_a) + Var(AUC\_b)
%     - 2 * Cov(AUC\_a, AUC\_b)).
%   - Two-sided p-value from standard normal.
%   - 15 pairwise comparisons total (C(6, 2) = 15 model pairs).
%   - No Bonferroni correction applied; acknowledged as a limitation in
%     Section 5.4.
%
% Earlier issue disclosure (T-R15): a prior version used only seed=0,
% making AUC and z-statistics depend on which initialisation happened to
% be labelled seed 0. Averaging across seeds before pooling removes this
% dependence.
%
% Cite: \cite{lopez_de_prado_2018} Chapters 11 and 14; \cite{delong_1988}.


\subsubsection{Stability diagnostics, model comparison, and ranking}\label{subsubsec:stability_ranking}

% ── Stability diagnostics ────────────────────────────────────────
%
% Topics:
%   - Feature stability: per-feature selection frequency across 28
%     folds. Frequency > 0.80 -> "stable". Flat profile -> diffuse signal
%     across many features (a finding, not a bug).
%   - FFD stability: mean and std of d* for ATR across 28 folds.
%     std(d*) < 0.1 -> consistent stationarity structure across periods;
%     > 0.1 -> time-varying persistence.
%
% ── Model comparison and ranking ─────────────────────────────────
%
% Topics:
%   - Primary criterion: median path Sharpe across 7 paths (descending).
%   - Tiebreaker: std Sharpe (ascending; prefer consistency).
%   - Columns: rank, model, median\_sharpe, std\_sharpe, DSR, mean\_F1,
%     mean\_acc, mean\_log\_loss, mean\_AUC, median\_max\_dd,
%     median\_cum\_ret, median\_win\_rate, median\_profit\_factor.
%   - Self-contained caption (Cochrane).


\subsection{Symbolic Extraction}\label{subsec:symbolic}

% Frame as exploratory analysis (per advisor guidance), not the core
% contribution. Keep main text approximately 1.5 pages. Detailed training
% logs, edge-by-edge R^2 values, unsimplified formulas -> Appendix F.


\subsubsection{Purpose, fold selection, and data preparation}\label{subsubsec:symbolic_setup}

% Invariance note (v6, 5 seeds). The symbolic-extraction outputs are
% identical to the previous 3-seed locked run. PyKAN is retrained
% from scratch on the same training fold (split 27) using a fixed
% internal seed, so the closed-form formula, the architecture
% [3, 3, 2], the symbolification rate 100\%, the pre-symbolic
% accuracy 46.67\%, the post-symbolic accuracy 51.85\%, and the
% phase-by-phase validation accuracies are all invariant under the
% outer-loop seed-count change. What does shift slightly is the
% locked-run KAN F1 on split 27 (the CPCV-evaluated efficient-kan
% number reported by the symbolic-extraction logger), which moves
% from 0.3048 in the previous run to 0.3232 in the current run
% because the F1 is averaged over the 5 outer seeds. The symbolic
% representation describes the same PyKAN model in both runs.
%
% Topics:
%   - Operates after all 28 CPCV splits have been evaluated.
%   - Reconstructs from prep\_info: stored d*, fitted scaler, selected
%     feature list. No per-instance state passed between efficient-kan
%     and PyKAN.
%   - Inherits architecture defaults from kan\_model.py: KAN\_HIDDEN=5,
%     KAN\_GRID=5, KAN\_K=3. Applies data-aware safety floor
%     (PYKAN\_MIN\_SAMPLES\_PER\_PARAM=5).
%
% ── Fold selection ───────────────────────────────────────────────
%
%   - fold\_selection="last" (split 27): test set covers most recent
%     data; closest to deployment scenario.
%   - Justification: most recent data reflects current regime.
%   - Acknowledged limitation: formula is fold-specific.
%   - Alternative "best" (highest KAN F1) maximises meaningful symbolic
%     output but biases toward best-performing fold.
%
% ── Data preparation ─────────────────────────────────────────────
%
%   - Re-generates CPCV splits for training indices.
%   - Applies stored d* to full series; extracts training fold.
%   - Applies stored RobustScaler.
%   - Selects features via the n\_top\_features parameter and the
%     feature\_selection\_strategy parameter. The current locked
%     configuration uses n\_top\_features=3 and
%     feature\_selection\_strategy=\"stability\" (ranks by 28-fold MDA
%     selection frequency, keeps the top 3). Earlier locked
%     configurations used n\_top\_features=5 with the same strategy.
%     The parameter trades input dimensionality (more features ->
%     more potential signal) against formula readability (fewer
%     features -> more compact closed form).
%   - 80/20 chronological split into model-train + validation.
%   - Tanh normalisation fitted on training portion:
%     z = tanh((x - mean) / (std + 1e-8)).
%   - Defensive input handling: prepare\_extraction\_data coerces y to
%     a pd.Series indexed on X.index regardless of how the caller
%     passes it (Series, numpy array, other array-like); length
%     mismatch raises a clear ValueError. Catches a common notebook
%     pattern where y gets shadowed by a pooled-prediction array.


\subsubsection{Four-step PyKAN algorithm}\label{subsubsec:pykan_algorithm}

% The four steps of \cite{cho_lee_kim_2025}'s Algorithm 1: train, prune,
% symbolify, affine fine-tune. All four steps share the same PyKAN
% model object and operate sequentially on the training fold's KAN.
%
% ── Step 1. Three-phase training ─────────────────────────────────
%
% Phases:
%   - Phase 1. Adam (lr=1e-3, wd=1e-3), 600 steps. Gaussian noise
%     injection (std=0.05) clamped to [-1, 1]; dropout-like regulariser;
%     early stopping on val loss.
%   - Phase 2a. LBFGS warmup (lr=0.01), 20 steps. No regularisation;
%     light refinement.
%   - Phase 2b. LBFGS sparsity (lr=0.01), 20 steps. L1 + entropy
%     regularisation (lamb=0.002); encourages sparse activations.
%
% Topics:
%   - Grid extension disabled (PYKAN\_GRID\_EXTEND=False): with 540
%     training samples and 135 validation samples (training fold of
%     split 27 after the 80/20 split, locked run), refining grid 3 to 5
%     adds parameters faster than data supports.
%   - Dynamic majority-baseline accuracy gate: after Phase 1 (Adam),
%     the pipeline checks val\_acc against
%     max(PYKAN\_MIN\_ACCURACY=0.53, val\_majority\_baseline + 0.01).
%     The gate adjusts to the actual class balance in the validation
%     fold rather than assuming the dataset-level base rate. The
%     previous configuration used a fixed val\_acc >= 0.53 threshold;
%     the dynamic gate is a methodology improvement that surfaces
%     sub-baseline learning more reliably. If the gate fails, log a
%     warning but continue. A second, structurally similar gate runs
%     at the pre-prune stage using the same dynamic formula.
%     Symbolic extraction may yield constants when the underlying
%     model has not learned signal.
%
% Architecture matches the CPCV-benchmarked KAN:
%   - The PyKAN model trained here uses the same single-hidden-layer
%     topology and the same width cap (width1 <= 6) as the efficient-kan
%     model evaluated under CPCV.
%   - The extracted symbolic formula therefore describes the same
%     architecture the benchmark numbers describe, not a simplified
%     surrogate.
%   - This is a structural change from earlier configurations, where the
%     symbolic-extraction KAN was constrained more tightly than the
%     CPCV-evaluated KAN to keep sympy simplification tractable; the
%     previous architecture-mismatch caveat
%     (PYKAN\_SYMBOLIC\_WIDTH\_CAP, PYKAN\_SYMBOLIC\_DROP\_WIDTH2,
%     PYKAN\_SYMBOLIC\_FORCE\_GRID) is no longer needed.
%
% ── Step 2. Pruning ──────────────────────────────────────────────
%
% Topics:
%   - Forward pass populates cached activations.
%   - model.attribute() for importance scoring.
%   - model.prune(threshold=PRUNE\_THRESHOLD=0.01) with multi-API
%     fallback (PyKAN versions vary).
%   - Verify pruned model still forward-passes; revert if broken.
%   - Typical compression with width1 <= 6: a [n\_features, 6, 2]
%     topology often prunes to something like [n\_features, 3, 2] after
%     training-time activation thresholds. Earlier configurations with
%     width1 up to 12 produced richer pre-prune networks but created the
%     architectural-mismatch problem described in Step 1 above.
%
% ── Step 3. Symbolification ──────────────────────────────────────
%
% Topics:
%   - Library (14 functions): x, x^2, x^3, x^4, exp, log, sqrt, tanh,
%     sin, cos, abs, sgn, arctan, 0.
%   - Per-edge: model.suggest\_symbolic(l, i, j, topk=5,
%     lib=SYMBOLIC\_LIBRARY) returns ranked candidates.
%   - Three format handlers (PyKAN versions return DataFrame, flat
%     tuple, or nested tuples).
%
% Constant-skip logic / brute-force fallback:
%   - When "0" wins (always wins total\_loss due to zero complexity
%     penalty), iterate manually through 14 candidates, call
%     fix\_symbolic for each, capture R^2 from PyKAN's stdout via regex
%     r"r2 is ([\textbackslash d.eE+-]+)", keep best non-constant
%     function. Restore model state between trials.
%
% Selection: best non-constant with R^2 >= SYMBOLIC\_R2\_THRESHOLD=0.30
% (lowered from 0.50; 0.50 was too aggressive in the weak-signal regime).
%
% Below threshold -> keep spline (partial symbolification).
%
% ── Step 4. Affine fine-tuning ───────────────────────────────────
%
% Topics:
%   - 30 LBFGS steps at lr=0.0004 (from VIX KAN paper) on remaining
%     affine parameters a, b, c, d per symbolic edge.
%   - Adjusts the symbolic fit without changing the function family.
%   - NaN detection reverts to pre-fine-tune state.


\subsubsection{Formula extraction and output}\label{subsubsec:formula_extraction}

% Topics:
%   - Tries three SymPy variable naming conventions: x1..xn, x\_1..x\_n,
%     x\_0..x\_{n-1}.
%   - model.symbolic\_formula(var=...) with no-var fallback.
%   - Two outputs: logit\_bearish (class 0), logit\_bullish (class 1).
%   - Decision function: logit\_bullish - logit\_bearish.
%   - P(up) = sigmoid(decision) = 1 / (1 + exp(-decision)).
%   - Substitute placeholder variables with feature names.
%   - sympy.simplify with 30-second timeout via threading (sympy hangs
%     on complex expressions).
%   - sympy.nsimplify(tolerance=1e-3) for cleaner rational coefficients.
%   - Identifies surviving features (those in decision.free\_symbols).
%
% Output dictionary contents:
%   - logit\_bearish, logit\_bullish, decision\_function, p\_up\_formula
%     (strings).
%   - sympy\_objects (SymPy expression objects for partial-derivative
%     sensitivity).
%   - pre\_symbolic\_accuracy, post\_symbolic\_accuracy.
%   - symbolification\_rate (fraction of edges symbolified).
%   - pruned\_architecture (e.g., [5, 3, 2]).
%   - surviving\_features.


\subsubsection{Known fragilities and acknowledged limitations}\label{subsubsec:symbolic_fragilities_limits}

% T-R15 (transparency).
%
% ── PyKAN fragilities encountered ────────────────────────────────
%
% Topics:
%   - 'sigmoid' is NOT in PyKAN's internal SYMBOLIC\_LIB -> KeyError.
%   - '1/x' causes division-by-zero at affine fine-tuning.
%   - 'abs' (the absolute value primitive) symbolises cleanly via
%     fix\_symbolic but yields a SymPy expression that is non-
%     differentiable at zero; for features whose dataset mean is close
%     to zero in the tanh-normalised space, partial-derivative
%     sensitivity at the mean returns NaN (Section 4.6, L6 in
%     Section 5.4).
%   - PyKAN uses 1-based variable naming (x\_1..x\_n).
%   - suggest\_symbolic return format varies across PyKAN versions.
%   - sympy.simplify can hang indefinitely (-> 30s timeout).
%   - "0" always wins total\_loss (-> brute-force fallback).
%
% ── Acknowledged limitations (also discussed in 5.3) ─────────────
%
% Topics:
%   - Fold-specific (different folds -> potentially different formulas).
%   - Small sample (540 training observations and 135 validation
%     observations after 80/20 split on split 27).
%   - R^2 threshold sensitivity (lowered to 0.30 to admit symbolic fits
%     in weak-signal regime).
%   - Post-symbolic accuracy may be lower than pre-symbolic accuracy.
%   - Brute-force fallback depends on PyKAN's stdout format.
%   - Architecture parity (no longer a limitation, formerly a caveat):
%     the single-hidden-layer / width1 <= 6 configuration now matches
%     the CPCV-benchmarked KAN, so the extracted formula describes the
%     same architecture the benchmark numbers describe.
%
% Cite: \cite{cho_lee_kim_2025}, \cite{liu_kan_2024},
% \cite{noorizadegan_2026}.




% ==================================================================
% CHAPTER 4: RESULTS
% ==================================================================
\newpage
\section{Results}\label{sec:results}

% Cochrane: lead with the main result. T-R10: frame empirics as theory
% tests. Numbers below come from the locked end-to-end run.


\subsection{Headline}\label{subsec:results_headline}

% T-R1, T-R3.
%
% Methodology note (v6): the locked configuration now uses 5 seeds
% across all six models (up from 3), producing 840 = 28 splits x 6
% models x 5 seeds prediction entries. Data alignment unchanged
% (1,168 events, 56.85\%/43.15\% class balance).
%
% Five opening sentences, each leading with the strongest claim:
%   - "No model achieves DSR >= 0.95. The top DSR among trained models
%      (KAN) is 0.2470, far below significance after correcting for
%      selection bias across the six compared models. Buy-and-hold
%      posts a higher median Sharpe (2.2722) than every trained model,
%      illustrating how punishing the AFML correction layer is for
%      trading strategies on a trending asset."
%   - "PBO = 0.657 lands in the adversarial regime (PBO > 0.5): in 23
%      of 35 IS/OOS partitions the in-sample best model underperforms
%      the out-of-sample median. The IS-winner is the OOS-loser more
%      often than not, which means model selection across this
%      six-model cross-section is not just unreliable but actively
%      misleading on more than half the sub-path partitions.
%      Leave-one-out PBO sharpens the picture: excluding logistic
%      reduces PBO most (Delta -0.286 to 0.371, back in the moderate-
%      overfitting band), followed by ar\_logistic and lstm (both
%      Delta -0.200 to 0.457), xgboost and kan (both Delta -0.114 to
%      0.543), with random\_forest PBO-neutral (Delta +0.000). Random
%      Forest is the structural finding: removing it does not change
%      the cross-section's IS/OOS rotation, so RF contributes nothing
%      to the adversarial reading. The previous run's 'no anchor'
%      pattern (every exclusion reduced PBO) does not reproduce in the
%      larger 5-seed configuration; instead one model anchors PBO from
%      below (RF) and the rest are destabilisers."
%   - "The previous run's KAN strict-positive bootstrap CI does NOT
%      reproduce. Under the uniform 5-seed configuration, 6 of 7
%      bootstrap 95\% CIs on the median Sharpe cross zero; only
%      buy-and-hold's CI (1.8903, 2.4685) stays strictly positive.
%      KAN's CI moves from the v5 strict-positive (0.0971, 1.2413) to
%      a zero-crossing (-0.4677, 1.1607) despite its median Sharpe
%      rising from 0.4778 to 0.5588 and its DSR rising from 0.2198 to
%      0.2470. The increased seed cardinality reduces within-model
%      averaging variance enough to widen the bootstrap interval on
%      the lower side; the headline KAN result that v5 reported is
%      therefore not robust to the seed-count change."
%   - "2 of 15 DeLong pairwise AUC tests reach alpha=0.05 in the 5-seed
%      configuration, the first time any AUC pair has cleared
%      significance in any locked run. AR Logistic vs XGBoost
%      (z=2.0263, p=0.0427) and Random Forest vs XGBoost (z=2.1466,
%      p=0.0318) are both significant in the direction of XGBoost
%      being the lower-AUC half (XGBoost's pooled AUC sits at 0.4942,
%      below the 0.5 baseline). The two near-significant pairs at
%      p ~ 0.054 in v5 (with 3 seeds) clear the threshold once the
%      within-model averaging is over 5 seeds. The structural reading:
%      XGBoost's pooled AUC is now distinguishable from the two
%      best-AUC trained models (AR Logistic, Random Forest) at
%      alpha=0.05, but no other pairwise difference clears, and 13 of
%      15 pairs remain indistinguishable from sampling noise."
%   - "The symbolic extraction pipeline produces a closed-form decision
%      function with 100\% symbolification rate (15 of 15 edges) on
%      three surviving features (eth\_btc\_ratio, log\_returns\_lag6,
%      natgas\_ret\_30). Post-symbolic accuracy is 51.85\%, a 5.18
%      percentage-point gain over the pre-symbolic 46.67\%, but both
%      numbers sit at or below the 50\% baseline: the PyKAN model
%      retrained on the deployment fold did not extract signal that
%      beats coin-flip. The symbolic-extraction outputs are identical
%      to the previous run because PyKAN is retrained with a fixed
%      seed on the same training fold (split 27 with
%      n\_top\_features=3 and feature\_selection\_strategy='stability'),
%      making the symbolic representation invariant under the
%      outer-loop seed count. The Phase-1 dynamic gate (max(0.53,
%      base\_rate + 0.01) = 0.5693) fails at val\_acc 0.4815, the
%      pre-prune gate fails at val\_acc 0.4667, and recovery in the
%      LBFGS phases does not get past 50\%. The Abs(.) non-
%      differentiability flag from earlier configurations is resolved:
%      the formula uses sin, cos, tanh, x\^2, x\^3, and x\^4 as
%      primitives, all smooth at the dataset mean."
%
% Roadmap: 4.2 model comparison, 4.3 classification, 4.4 financial,
% 4.5 stability, 4.6 symbolic.


\subsection{Model Comparison}\label{subsec:results_comparison}

% This is the "main result" per Cochrane.
%
% Self-contained caption (Table X):
%   "Model comparison ranked by median Sharpe over seven CPCV backtest
%   paths, with buy-and-hold included as a naive baseline. DSR is
%   computed against n\_trials=6 (number of compared models), correcting
%   the observed Sharpe for selection bias and non-normal returns; DSR
%   is not defined for buy-and-hold since it is not part of the model
%   selection pool. The Sharpe confidence interval is a bootstrap 95\%
%   CI on the per-path median Sharpe. No model achieves DSR >= 0.95.
%   Tiebreaker: standard deviation of path Sharpe ascending."
%
% Columns: Rank | Model | Median Sharpe | Sharpe CI (low, high) |
%          Std Sharpe | DSR | Med Sortino | Med Calmar | Mean F1 |
%          Mean Acc | Mean AUC | Med Max DD | Med Cum Ret |
%          Med Win Rate | Med Profit Factor
%
% Locked-run v6 numbers (5 seeds across all six models; paste into
% the table):
%   1 | Buy-and-Hold          |  2.2722 | ( 1.8903, 2.4685) | 0.5768 | n/a    |  2.3390 |  2.4111 | n/a    | n/a    | n/a    | -0.9998 | 45896.39 | 0.5559 | 1.3007
%   2 | KAN                   |  0.5588 | (-0.4677, 1.1607) | 0.7539 | 0.2470 |  0.3211 |  0.0820 | 0.4070 | 0.5318 | 0.5103 | -0.2466 |   0.0841 | 0.5315 | 1.1290
%   3 | AR Logistic           |  0.3806 | (-0.0544, 1.2101) | 0.7355 | 0.1896 |  0.3894 |  0.0408 | 0.4666 | 0.5240 | 0.5204 | -0.2739 |   0.0852 | 0.5318 | 1.0687
%   4 | LSTM                  |  0.1541 | (-0.3902, 0.7032) | 0.6388 | 0.1291 |  0.1189 |  0.0075 | 0.4238 | 0.5179 | 0.5105 | -0.2986 |   0.0102 | 0.5109 | 1.0293
%   5 | XGBoost               |  0.0493 | (-0.0638, 0.2514) | 0.5830 | 0.1065 |  0.0259 | -0.0024 | 0.4271 | 0.5185 | 0.5008 | -0.2172 |  -0.0027 | 0.5122 | 1.0125
%   6 | Random Forest         |  0.0287 | (-0.6165, 0.8834) | 0.7533 | 0.1023 |  0.0257 | -0.0190 | 0.4390 | 0.5248 | 0.5166 | -0.3775 |  -0.0579 | 0.5022 | 1.0109
%   7 | Logistic Regression   | -0.2173 | (-0.7321, 2.0594) | 1.5969 | 0.0618 | -0.1872 | -0.0353 | 0.4355 | 0.5241 | 0.5156 | -0.3874 |  -0.0797 | 0.5018 | 0.9703
%
% Methodology note (5-seed configuration). This run uses 5 seeds
% across all six models (up from 3 in the previous configuration),
% producing 840 = 28 splits x 6 models x 5 seeds prediction entries.
% The within-model averaging is now over 5 independent inner-CV
% trajectories per (model, split) cell; the path-level Sharpe matrix
% that feeds DSR and PBO is the same 7 x 6 shape, but every entry is
% now averaged over 5 seed-replicates rather than 3. The seed-count
% change is the only methodology difference between this run and the
% previous one; the data alignment (1,168 events, 56.85\%/43.15\%
% class balance, 2015-08-08 to 2026-05-02), the feature pool (73
% features), and the symbolic-extraction call (n\_top\_features=3,
% feature\_selection\_strategy="stability", fold 27) are unchanged.
%
% Key observations (state facts here, save interpretation for 5.1):
%   - DSR. Top DSR among trained models is KAN at 0.2470 (up from
%     0.2198 in the previous run); all six trained-model values fall
%     below 0.95. No model demonstrates predictive ability that
%     survives correction for selection bias. The top three DSRs
%     (KAN 0.2470, AR Logistic 0.1896, LSTM 0.1291) span 0.12, with
%     smaller gaps to XGBoost (0.1065), Random Forest (0.1023), and
%     Logistic Regression (0.0618). Buy-and-hold is excluded from the
%     DSR computation because it is not an Optuna-tuned model and does
%     not enter the selection pool. KAN's DSR rose by 0.027 between
%     the two runs; Random Forest's DSR fell from 0.1415 to 0.1023 (a
%     0.04 drop) and is the largest negative change in the
%     cross-section.
%   - Sharpe confidence intervals (key finding REVERSED from previous
%     run). 6 of 7 entries have a 95\% bootstrap CI that crosses zero;
%     only buy-and-hold (1.8903, 2.4685) stays strictly positive. The
%     previous run's KAN strict-positive finding (CI 0.0971, 1.2413)
%     does NOT reproduce in the 5-seed configuration: KAN's CI moves
%     to (-0.4677, 1.1607), which crosses zero on the lower side.
%     This shift happens despite KAN's median Sharpe rising from
%     0.4778 to 0.5588 and its DSR rising from 0.2198 to 0.2470. The
%     mechanism is the within-model variance change: averaging the
%     path-Sharpe matrix over 5 seeds rather than 3 widens the
%     bootstrap CI's lower tail more than it tightens the median; the
%     headline KAN evidence from the previous run is therefore not
%     robust to the seed-count change.
%   - Buy-and-hold reference point. Buy-and-hold posts the highest
%     median Sharpe (2.2722) and the highest cumulative return
%     (approximately 45,896x over the 2015 to 2026 period, reflecting
%     BTC's trajectory from approximately \$325 to roughly \$108k). Its
%     median max drawdown is -0.9998 (essentially -100\%): the position
%     experiences the full BTC drawdown profile, which is precisely
%     what bet sizing and the dual-weighted loss are designed to
%     mitigate. Trained models trade off return for downside
%     protection: KAN's cumulative return of 0.0841 is roughly six
%     orders of magnitude smaller than buy-and-hold's, but its max
%     drawdown of -0.2466 is four times shallower. This is the
%     financial-management framing in which the DSR result should be
%     read.
%   - Ranking. KAN leads the trained models with median Sharpe 0.5588
%     (up from 0.4778 in the previous run); AR Logistic third at
%     0.3806 (unchanged); LSTM fourth at 0.1541 (down from 0.1808);
%     XGBoost fifth at 0.0493 (~= 0.0484); Random Forest sixth at
%     0.0287 (down from 0.2047); Logistic Regression last at -0.2173
%     (the only model with a negative median Sharpe, slightly worse
%     than the previous run's -0.1545). The largest single shifts are
%     KAN moving up by 0.081 in median Sharpe and Random Forest moving
%     down by 0.176. The previous run's top-3 trained models were
%     \{KAN, AR Logistic, Random Forest\}; the current top-3 are
%     \{KAN, AR Logistic, LSTM\}, with Random Forest dropping out and
%     LSTM moving up.
%   - AUC compression. All mean AUCs sit in the 0.5008 to 0.5204 band.
%     Models are nearly indistinguishable at the classification level;
%     bet sizing and the dual-weighted loss are what translate the
%     tiny edges into the wider Sharpe spread. AR Logistic posts the
%     highest mean AUC (0.5204), narrowly ahead of Random Forest
%     (0.5166), Logistic Regression (0.5156), LSTM (0.5105), KAN
%     (0.5103), and XGBoost (0.5008). The pooled-AUC ordering is
%     similar to the previous run but two pairwise differences now
%     reach DeLong significance at alpha=0.05 (Section 4.3).
%   - Q1, Q2 evidence (feature engineering). AR Logistic again posts
%     a positive median Sharpe (0.3806) and the highest mean AUC;
%     pure-momentum lags continue to deliver competitive financial
%     performance and the highest AUC of the cross-section. Logistic
%     Regression is still last on median Sharpe (-0.2173). Random
%     Forest drops to sixth place this run, the largest mover in the
%     ranking; its DSR fell by 0.04 and its median Sharpe fell from
%     0.2047 to 0.0287. The deeper sample within each model (5 seeds
%     x 28 splits = 140 path-level Sharpe estimates per model
%     entering the median computation, vs. 84 in the previous run)
%     penalises models whose path-Sharpe distribution is skewed by a
%     small number of strong paths.
%   - Q3 evidence (KAN positioning). KAN is the top trained model on
%     median Sharpe (0.5588) and on DSR (0.2470) for the second
%     consecutive run. Mean accuracy (0.5318) is mid-pack, mean AUC
%     (0.5103) is competitive but not first, median win rate 0.5315
%     is third among trained models (AR Logistic 0.5318, XGBoost
%     0.5122). Median max-drawdown -0.2466 is the best (smallest)
%     among trained models. The previous run's strict-positive
%     bootstrap-CI finding does not reproduce; the headline KAN
%     evidence is now "rank-1 trained Sharpe and rank-1 trained DSR
%     over two consecutive runs with different seed cardinalities,
%     but the bootstrap CI on the median Sharpe crosses zero". The
%     leave-one-out PBO finding in Section 4.4 reads differently from
%     both the v4 KAN-anchor pattern and the v5 no-anchor pattern:
%     KAN's exclusion in v6 drops PBO from 0.657 to 0.543 (Delta
%     -0.114), the second-smallest destabilising contribution in the
%     LOO table, while Random Forest is PBO-neutral.
%   - Std Sharpe outlier. XGBoost now posts the lowest std Sharpe
%     (0.5830) among trained models, with a correspondingly narrow CI
%     (-0.0638, 0.2514) that crosses zero by a small margin. LSTM's
%     std (0.6388) and AR Logistic's std (0.7355) are the
%     next-tightest. Logistic Regression carries the highest std
%     Sharpe (1.5969) and the widest CI (-0.7321, 2.0594), spanning
%     nearly three Sharpe units. The wide CI on Logistic Regression
%     contributes to its top destabilising role in the LOO PBO (Delta
%     -0.286 when excluded).


\subsection{Classification Metrics}\label{subsec:results_classification}

% Topics:
%   - Per-split F1 distribution. Mean F1 ranges narrowly from 0.4070
%     (KAN) to 0.4666 (AR Logistic). The high-F1 / competitive-Sharpe
%     combination for AR Logistic continues to confirm that
%     classification quality and bet-sizing translation are partially
%     independent: pure momentum gets the average call right slightly
%     more often than the engineered-feature models, and its
%     bet-sizing translation now produces a competitive third-place
%     median Sharpe in the 5-seed configuration.
%   - Pooled AUC. Per-model pooled AUC across all 28 splits sits in
%     the 0.4942 to 0.5149 band. All models hover within two
%     percentage points of random; the entire signal contest plays
%     out inside an AUC range smaller than what most reviewers would
%     consider economically material. XGBoost's pooled AUC (0.4942)
%     sits below 0.5 for the second consecutive run; in the 5-seed
%     configuration this below-chance positioning is now significant
%     against AR Logistic and Random Forest (DeLong table below).
%
% DeLong pairwise AUC: 2 of 15 pairs significantly different at
% alpha=0.05 (the first locked configuration to clear DeLong
% significance on any pair). Pooled across splits, averaged across
% the 5 seeds per model.
%
% Locked-run v6 pairs (paste into a table):
%   Pair                            | AUC\_a | AUC\_b | Delta AUC | z       | p      | Significant
%   AR Logistic vs Logistic         | 0.5117 | 0.5084 | +0.0032   |  0.3982 | 0.6905 | No
%   AR Logistic vs Random Forest    | 0.5117 | 0.5068 | +0.0048   |  0.5671 | 0.5706 | No
%   AR Logistic vs XGBoost          | 0.5117 | 0.4942 | +0.0175   |  2.0263 | 0.0427 | YES
%   AR Logistic vs LSTM             | 0.5149 | 0.5048 | +0.0100   |  0.9918 | 0.3213 | No
%   AR Logistic vs KAN              | 0.5117 | 0.5062 | +0.0055   |  0.7018 | 0.4828 | No
%   Logistic vs Random Forest       | 0.5084 | 0.5068 | +0.0016   |  0.2528 | 0.8005 | No
%   Logistic vs XGBoost             | 0.5084 | 0.4942 | +0.0143   |  1.7728 | 0.0763 | No
%   Logistic vs LSTM                | 0.5100 | 0.5048 | +0.0052   |  0.4964 | 0.6196 | No
%   Logistic vs KAN                 | 0.5084 | 0.5062 | +0.0022   |  0.3596 | 0.7191 | No
%   Random Forest vs XGBoost        | 0.5068 | 0.4942 | +0.0127   |  2.1466 | 0.0318 | YES
%   Random Forest vs LSTM           | 0.5091 | 0.5048 | +0.0043   |  0.4331 | 0.6650 | No
%   Random Forest vs KAN            | 0.5068 | 0.5062 | +0.0006   |  0.0932 | 0.9258 | No
%   XGBoost vs LSTM                 | 0.4956 | 0.5048 | -0.0092   | -0.9208 | 0.3572 | No
%   XGBoost vs KAN                  | 0.4942 | 0.5062 | -0.0121   | -1.6872 | 0.0916 | No
%   LSTM vs KAN                     | 0.5067 | 0.5073 | -0.0006   | -0.0561 | 0.9552 | No
%
% Pairwise summary:
%   - Two pairs clear alpha = 0.05. Random Forest vs XGBoost at
%     z = 2.1466 (p = 0.0318) and AR Logistic vs XGBoost at z = 2.0263
%     (p = 0.0427); both significant in the direction of XGBoost
%     being the lower-AUC half. The two near-significant pairs at
%     p ~ 0.054 in the previous 3-seed run clear the threshold once
%     the within-model averaging is over 5 seeds. The structural
%     reading: XGBoost's pooled AUC (0.4942) is now distinguishable
%     from the two best-AUC trained models (AR Logistic 0.5117 to
%     0.5149 and Random Forest 0.5068 to 0.5091) at alpha = 0.05.
%   - The 13 non-significant pairs include all KAN comparisons (p in
%     the 0.09 to 0.96 range against the five other trained models)
%     and all LSTM and Logistic pairwise comparisons. The KAN-vs-
%     XGBoost pair is the closest non-significant pair on the KAN
%     side at z = -1.6872 (p = 0.0916).
%   - This is the first locked configuration in which any DeLong
%     pair clears significance. The previous three runs all returned
%     0 of 15. The 2-of-15 finding is a genuine new positive result
%     from the 5-seed increase: doubling the within-model averaging
%     reduces the noise floor enough for the largest AUC differences
%     in the cross-section to become statistically distinguishable.
%
% Effect-size disclosure (T-R15: write for reviewers):
%   - The entire AUC range across all six models is 0.021 (from
%     XGBoost at 0.4942 to AR Logistic at 0.5149), the same order of
%     magnitude as the seed-to-seed variation. The two significant
%     pairs both lean on the lower-AUC tail of the cross-section
%     (XGBoost as the consistent lower half), so the practical
%     reading is "XGBoost's AUC is significantly lower than the top
%     two trained models", not "any pair of trained models has
%     significantly different classification accuracy from any
%     other".
%   - No Bonferroni correction is applied. With 15 comparisons and a
%     family-wise alpha = 0.05 / 15 ~= 0.0033, neither significant
%     pair survives Bonferroni. The 0.0318 and 0.0427 p-values do
%     clear Benjamini-Hochberg at q = 0.05 if applied to the 15-pair
%     family (the Random Forest vs XGBoost pair clears BH(2/15) =
%     0.0067, and AR Logistic vs XGBoost clears BH(4/15) = 0.0133),
%     so the finding is robust to the less-strict FDR correction
%     but not to Bonferroni.
%
% Confusion matrices (compact, Appendix D for full): aggregated
% TP/FP/TN/FN per model. Most models post a slight class-1 bias
% consistent with the empirical 56.85\% / 43.15\% Up/Down base rate.


\subsection{Financial Performance}\label{subsec:results_financial}

% Topics:
%   - Path-level Sharpe distribution (v6, 5 seeds). Per-model row x 7
%     paths + median + std. KAN posts the highest trained-model
%     median (0.5588) but its std Sharpe (0.7539) gives a bootstrap
%     95\% CI (-0.4677, 1.1607) that crosses zero on the lower side;
%     the previous run's strict-positive KAN CI does not reproduce.
%     AR Logistic (median 0.3806, std 0.7355) sits second with a CI
%     (-0.0544, 1.2101) that crosses zero very narrowly on the lower
%     side. LSTM (median 0.1541, std 0.6388) is third, CI (-0.3902,
%     0.7032). XGBoost (median 0.0493, std 0.5830) is fourth with
%     the narrowest CI (-0.0638, 0.2514) of any trained model.
%     Random Forest (median 0.0287, std 0.7533) is fifth, CI
%     (-0.6165, 0.8834). Logistic Regression (median -0.2173, std
%     1.5969) is the worst trained model with the widest CI
%     (-0.7321, 2.0594), spanning nearly three Sharpe units in both
%     directions. Buy-and-hold (median 2.2722, std 0.5768) dominates
%     the median ranking with the tightest CI (1.8903, 2.4685) of
%     the seven entries and is the only entry whose CI stays
%     strictly positive.
%   - Sortino and Calmar. The downside-deviation-adjusted Sortino
%     ratios reorder the trained models: AR Logistic (0.3894), KAN
%     (0.3211), LSTM (0.1189), XGBoost (0.0259), Random Forest
%     (0.0257), Logistic Regression (-0.1872). AR Logistic's Sortino
%     leads the trained models despite KAN's higher median Sharpe,
%     suggesting AR Logistic's path-Sharpe distribution loads more
%     variance on the upside than on the downside. KAN's Sortino is
%     second. Calmar (Sharpe over max drawdown) compresses the
%     trained models into a narrow -0.04 to +0.08 band, with KAN at
%     0.0820 leading among trained models. Buy-and-hold's Calmar
%     (2.4111) is two orders of magnitude above any trained model,
%     but its near-100\% max drawdown means a single rebalancing
%     decision during the deepest crash would wipe out the strategy.
%   - DSR computation detail. For KAN (top-ranked trained model in
%     this run): observed median Sharpe 0.5588 over 7 paths, DSR =
%     0.2470 against n\_trials=6. The expected maximum Sharpe under
%     the null with six compared models, combined with the pooled
%     skew/kurt and SE(SR), produces a DSR threshold the observed
%     Sharpe does not clear. The 0.2470 figure says KAN's median
%     Sharpe is roughly at the 25th percentile of what the null
%     hypothesis (all true Sharpes equal zero) would produce when
%     testing six independent strategies, far short of the 0.95
%     significance threshold. KAN's DSR rose from 0.2198 in the
%     previous run to 0.2470 here despite the bootstrap-CI evidence
%     becoming weaker: DSR depends on the median Sharpe and the
%     return-distribution skew/kurt, while the bootstrap CI depends
%     on the spread of the 7 path-Sharpes around the median; the two
%     diagnostics moved in opposite directions across the seed-count
%     change. Full DSR-component breakdown for the top model goes to
%     Appendix D.
%
% PBO result. PBO = 0.6571 = 23 of 35 IS/OOS partitions.
% Interpretation in context of DSR < 0.95:
%   - Theoretical regime map: PBO ~ 0 = robust selection (the IS-best
%     model tends to remain best out of sample); PBO ~ 0.5 = model
%     selection is random; PBO > 0.5 = adversarial.
%   - Observed regime: adversarial overfitting. In 65.71\% of
%     partitions the model that wins on three IS paths underperforms
%     the median of the four OOS paths. The IS ranking is
%     anti-informative in this run: a strategy that picks the IS-best
%     model and bets on it OOS does worse than the OOS median in
%     nearly two out of three partitions. This is consistent with
%     the previous run's adversarial finding (PBO 0.629 -> 0.657, a
%     0.028 increase) and represents the second consecutive locked
%     configuration in the adversarial regime. The PBO trajectory
%     across the four configurations evaluated (0.26 -> 0.37 -> 0.34
%     -> 0.629 -> 0.657) was stable in the moderate-overfitting band
%     for the first three runs and adversarial across the last two
%     runs; the cross-section's IS rankings are reproducibly
%     anti-informative under the 73-feature pool with uniform
%     seeds-per-model.
%   - Leave-one-out PBO reveals a structural finding (different from
%     the previous run's "no anchor" pattern). Re-running PBO on the
%     5x7 Sharpe matrix with one model dropped at a time produces a
%     pattern where Random Forest is PBO-neutral and the rest are
%     destabilisers. Sorted by 5-model PBO ascending:
%
%       Model excluded    5-model PBO    Delta vs baseline
%       ---------------------------------------------------
%       logistic              0.371          -0.286
%       ar_logistic           0.457          -0.200
%       lstm                  0.457          -0.200
%       xgboost               0.543          -0.114
%       kan                   0.543          -0.114
%       random_forest         0.657          +0.000
%
%   - Reading the LOO table. Random Forest is the anchor model in
%     this configuration: its exclusion leaves PBO unchanged at
%     0.657. This is structurally different from both (a) the v4
%     "KAN-anchor / one-model-anchor" pattern where one model's
%     presence stabilised the cross-section by being well-calibrated
%     and consistent across sub-paths, and (b) the v5 "no anchor"
%     pattern where every model's exclusion reduced PBO. In v6,
%     Random Forest is far enough from the IS-best slot across all
%     sub-path partitions that removing it does not change the IS/OOS
%     rotation among the remaining five models. Logistic Regression
%     is the largest destabiliser (Delta -0.286 to 0.371, back in
%     the moderate-overfitting band); excluding AR Logistic or LSTM
%     reduces PBO to 0.457 (Delta -0.200, just past the random-
%     selection threshold on the lower side); KAN and XGBoost are
%     equal destabilisers at Delta -0.114. The reading: the IS/OOS
%     rotation is concentrated in the top-five trained models, with
%     Random Forest sitting too low in the IS ranking to participate.
%     Two of the three locked-configuration runs with uniform
%     seeds-per-model now produce an anchor model under LOO PBO (KAN
%     in v4 from below, Random Forest in v6 from below); the v5
%     no-anchor pattern is now the outlier across configurations.
%   - Joint reading with DSR < 0.95: model selection within this
%     six-model universe is adversarial in the current configuration
%     (PBO 0.657 > 0.5), no model survives multiple-testing
%     correction (top DSR 0.2470 << 0.95), and the IS-best model is
%     the OOS-loser in two of three partitions. The previous run's
%     KAN strict-positive bootstrap CI does not reproduce in the
%     current run; 6 of 7 trained-model CIs cross zero, and only
%     buy-and-hold's CI stays strictly positive. The honest
%     conclusion is "no trained model has predictive ability that
%     survives multiple-testing correction; KAN is the top trained
%     model by median Sharpe and DSR but its bootstrap CI now
%     crosses zero; model selection in this cross-section is
%     structurally adversarial; Random Forest sits low enough in the
%     ranking to be PBO-neutral; and a naive buy-and-hold strategy
%     posts a 2.27 median Sharpe over the same period at the cost
%     of a -99.98\% maximum drawdown".
%
% Methodological implication (T-R15: write for reviewers, this is
% the AFML pay-off):
%   - Without PBO, the natural conclusion from the v6 model-comparison
%     table is "KAN achieved median Sharpe 0.56 on BTC daily direction
%     across 5 seeds, the top trained-model result of the
%     cross-section".
%   - With the headline PBO of 0.657 the conclusion sharpens: the
%     IS-best in any sub-path partition is the OOS-loser more often
%     than not, so even the top trained-model Sharpe must be read
%     against the structural risk that the IS lead does not carry to
%     OOS.
%   - The leave-one-out PBO refines the message further: Random
%     Forest is the only PBO-neutral model in the cross-section,
%     suggesting it sits below the IS/OOS rotation rather than
%     contributing to it.
%   - The PBO and DSR diagnostics are complementary: PBO answers
%     "did we pick the right model?", DSR answers "is the right
%     model good enough?"; here the answer is "no, the cross-section
%     is adversarial" to the first and "no" to the second.
%   - The DeLong 2-of-15 finding (Section 4.3) does not change the
%     PBO/DSR picture: significant AUC differences between AR
%     Logistic / Random Forest vs XGBoost establish that XGBoost is
%     the lower-AUC trained model, but no positive AUC difference
%     above 0.5 has been demonstrated, so the classification-side
%     significance does not propagate to a financial-side
%     significance claim.
%
% Equity curves figure: all 6 models + buy-and-hold overlay on the
% median (or best) path. Note specifically the XGBoost, Random
% Forest, and Logistic Regression equity curves, which are the three
% ending below 1.0 (median cumulative returns -0.0027, -0.0579, and
% -0.0797 respectively). Buy-and-hold's curve dominates the plot
% vertically (45,896x cumulative return); a log-scale y-axis is
% needed to see the trained models alongside it.
%
% Additional metrics table: cumulative return, annualised return,
% max DD, time under water, win rate, profit factor, mean |bet|,
% \% traded. Locked-run v6 highlights from the table at 4.2: AR
% Logistic median cumulative return 0.0852 (best among trained),
% KAN 0.0841, LSTM 0.0102, XGBoost -0.0027, Random Forest -0.0579,
% Logistic Regression -0.0797. Median max-drawdown ranges from
% -0.2172 (XGBoost, best among trained) to -0.3874 (Logistic
% Regression, worst); KAN's -0.2466 is the second-best. Buy-and-
% hold's median max-drawdown is -0.9998, the structural cost of
% holding through the 2017-18 and 2022 BTC crashes. KAN posts the
% second-best max-drawdown among trained models alongside its
% rank-1 trained median Sharpe; the combination is consistent with
% well-calibrated bet sizing protecting against the largest
% drawdowns.


\subsection{Feature Selection Stability and FFD Stability}\label{subsec:results_features}

% Topics:
%   - Feature selection frequency (per-feature count / 28 folds). Locked-
%     run profile is flat to a striking degree. Of 73 candidate features:
%       * 0 features stable (selected in more than 80\% of folds, > 22
%         of 28 folds).
%       * 2 features moderate (50\% to 80\%): eth\_btc\_ratio at 57.1\%
%         (16 of 28 folds), log\_returns\_lag6 at 53.6\% (15 of 28
%         folds).
%       * 71 features low (selected in fewer than 50\% of folds).
%       * 0 features never selected (the previous configuration had
%         tx\_count\_roc\_14 as the single never-selected feature; in
%         the current run no feature posts zero selection across the
%         28 folds).
%   - Headline finding. No compact feature subset consistently
%     dominates across time periods. The MDA-selected set turns over
%     substantially across the 28 folds: the median feature is selected
%     in roughly a third of folds, and even the most-stable feature
%     (eth\_btc\_ratio) clears 50\% by about 7 percentage points. This
%     is itself a finding: it suggests that the information relevant to
%     BTC daily direction is distributed across many features, with
%     regime-specific pockets of relevance, rather than concentrated in
%     a small permanent core. The stability profile shifts substantially
%     from the previous locked configuration: eth\_btc\_ratio moves
%     from 60.7\% to 57.1\% (small drop), bb\_width drops below the
%     50\% threshold (54\% -> low bucket), skewness also drops below
%     50\% (was 46\%, similar position), and log\_returns\_lag6 emerges
%     as the new second-most-stable feature at 53.6\% (the previous
%     run had log\_returns\_lag1 lower in the stability ranking but not
%     in the top-2 moderate set).
%   - Group breakdown of the two moderate features.
%       * 1 of 2 from the external / crypto-macro group
%         (eth\_btc\_ratio, the CoinMetrics-sourced ETH/BTC alt-rotation
%         signal).
%       * 1 of 2 from the lag / autoregressive group
%         (log\_returns\_lag6, the 6-day lag of log returns). First
%         time a lag feature appears in the cross-fold moderate-
%         stability set in any locked configuration.
%       * 0 from technical-analysis, macro, on-chain, or mathematical
%         (AFML Part 4) groups.
%   - On-chain question (Q2). No on-chain feature reaches 50\%
%     selection frequency. All eight on-chain columns (active
%     addresses ROC, transaction count ROC, hash rate ROC, MVRV, net
%     exchange flow, fee per transaction, exchange supply percent,
%     issuance) sit in the "low" bucket. The "free lunch hypothesis"
%     (on-chain transparency providing measurable signal beyond price
%     and volume) does not survive multi-model MDA + CPCV evaluation
%     in this dataset; on-chain features may carry information for
%     individual splits, but no on-chain column is consistently
%     selected across folds. Q2 answer: on-chain features do not
%     exhibit stability under the AFML feature-selection regime
%     applied here. This finding hardens across all four locked
%     configurations evaluated: in the 62-feature run
%     exchange\_supply\_pct survived to the symbolic extraction step
%     despite low cross-fold stability; in the previous 73-feature
%     2-seed run no on-chain feature reached the top-5 stable set;
%     in the previous 73-feature 3-seed run one on-chain feature
%     (tx\_count\_roc\_14) was never selected at all; in the current
%     run no on-chain feature is moderate.
%   - TA-feature drop. The previous locked configuration had bb\_width
%     and skewness (both TA / distribution-shape and volatility) in
%     the moderate / near-moderate set. In the current run, both drop
%     below the 50\% threshold. The TA group does not contribute any
%     moderate-stability feature; the moderate set is now {crypto-
%     macro, lag} rather than {crypto-macro, TA}.
%   - Figure suggestion: horizontal bar chart of selection frequency
%     (73 bars, sorted descending, coloured by feature group:
%     green=TA, blue=mathematical, orange=external macro, red=external
%     crypto-macro, brown=external on-chain, grey=lag).
%
% FFD stability. ATR is the only column subject to FFD (the only
% feature flagged non-stationary by the cross-fold ADF audit at
% alpha=0.05 with the FFD column whitelist ['atr']; the full ADF
% report identifies 8 non-stationary features at alpha=0.05, all
% retained without FFD because the locked configuration applies the
% transform only to ATR by design). Locked-run statistics across 28
% folds x 5 seeds (140 d* estimates):
%   - Mean d* = 0.198, std d* = 0.085.
%   - Range: d* in [0.050, 0.400]; modal value d* = 0.15 to 0.20.
%   - Std d* < 0.1 -> consistent stationarity structure across time
%     periods. The ATR series requires fractional differencing of
%     order ~ 0.20 regardless of the training-fold partition,
%     indicating that the persistence of the volatility-range process
%     is a stable property rather than a regime-conditional one. No
%     fold required d* = 1.0 (full integer differencing). The mean d*
%     shifted from 0.170 (previous run) to 0.198 here, a small upward
%     drift consistent with the slightly larger dataset (1,168 vs.
%     1,159 aligned events).
%
% Methodological note. Feature-selection turnover (high) and FFD-d*
% turnover (low) tell a coherent story: the noise-versus-signal ratio
% for individual features changes across regimes (only eth\_btc\_ratio
% and log\_returns\_lag6 clear 50\%), but the underlying time-series
% property that motivates fractional differencing is stable. Stability
% of d* validates the FFD-only-on-ATR design choice in 3.6.1;
% instability of feature selection validates the use of multi-model
% MDA over single-model SFI in 3.6.3.


\subsection{Symbolic Extraction Results}\label{subsec:results_symbolic}

% Q4 answer.
%
% ── Extraction summary ────────────────────────────────────────────
%
%   - Fold used: split 27 (groups 6 and 7), fold\_selection="last".
%     Test set covers the most recent CPCV partition; closest to
%     deployment scenario. Locked-run KAN F1 on this fold: 0.3232
%     (averaged over 5 seeds; previous 3-seed run posted 0.3048 on
%     the same split, so the F1 ticked slightly upward with the
%     additional seeds, but both numbers remain below KAN's
%     cross-split mean F1 of 0.4070 in this run).
%   - Methodology change: n\_top\_features reduced from 5 to 3. The
%     locked notebook call is run\_symbolic\_extraction(...
%     n\_top\_features=3, fold\_selection="last",
%     feature\_selection\_strategy="stability"). The strategy parameter
%     selects on cross-fold stability rather than per-fold MDA, and the
%     smaller feature count tightens the architecture to a [3, 3, 2]
%     topology that is humanly readable as a four-outer-term formula on
%     three inputs.
%   - Top-3 stable features entering extraction. eth\_btc\_ratio (57\%),
%     log\_returns\_lag6 (54\%), natgas\_ret\_30 (39\%). 70 features
%     excluded. The previous run extracted on five features including
%     bb\_width, skewness, and log\_returns\_lag1; the current top-3
%     set has zero overlap with the previous 4th and 5th features. The
%     two top-3 holdovers from the previous run are eth\_btc\_ratio
%     (\#1 in both) and natgas\_ret\_30 (\#5 -> \#3). The new entrant
%     is log\_returns\_lag6, displacing log\_returns\_lag1 as the
%     lag-group representative.
%   - Training sample size: 540 train + 135 val (after the 80/20 split
%     on the split-27 training fold, slightly larger than the previous
%     run's 530 + 133 due to the updated 1,168-sample alignment). 3
%     features in the input matrix.
%   - Architecture (data-aware fallback). The locked Optuna study did
%     not produce tuned KAN hyperparameters for split 27, so the
%     data-aware fallback fires. Final architecture [3, 3, 2]: 3 input
%     features -> 3-unit hidden layer -> 2-unit binary output. 15 active
%     edges, ~90 spline parameters for 540 training samples
%     (parameter-to-sample ratio 6.0x, comfortably above the 5x floor).
%     The hidden width (3) increased from 2 in the previous run,
%     reflecting the data-aware fallback's response to fewer input
%     features.
%
% ── Training diagnostics ─────────────────────────────────────────
%
%   - Phase 1 (Adam, 600 steps, weight decay 1e-3, noise std 0.05):
%     final train accuracy 0.5759, validation accuracy 0.4815.
%   - Phase-1 dynamic gate FAILS (val\_acc 0.4815 < 0.5693, where the
%     effective threshold is max(0.53, base\_rate + 0.01) = max(0.53,
%     0.5593 + 0.01) = 0.5693). The dynamic gate is a methodology
%     improvement: instead of a fixed 0.53 threshold, the gate now
%     adjusts to the actual class balance in the validation fold
%     (majority\_baseline + 0.01 margin, floored at
%     PYKAN\_MIN\_ACCURACY=0.53). The pipeline logs the warning "Adam
%     phase val\_acc=0.4815 < 0.5693 (majority baseline 0.5593 + 0.01
%     margin, floored at PYKAN\_MIN\_ACCURACY=0.53). PyKAN may not have
%     learned meaningful patterns" and continues to symbolification per
%     the design.
%   - Phase 2a (LBFGS warmup, 20 steps, no regularisation): train
%     0.6111, val 0.4296 (the only locked-run configuration where Phase
%     2a drops validation accuracy below Phase 1).
%   - Phase 2b (LBFGS sparsity, 20 steps, lambda=2e-3): train 0.6500,
%     val 0.4667.
%   - Pre-prune gate fails too. The pruning step's pre-prune check
%     (val\_acc < 0.5359, where 0.5359 = 0.5259 majority baseline + 0.01
%     margin) also fails at val\_acc 0.4667. The model has not learned
%     patterns that beat the majority baseline by even one percentage
%     point at any of the four training phases.
%   - Grid extension skipped (n\_train=540 <= 1000, grid stays at 3).
%   - Detailed phase-by-phase logs in Appendix F.
%
% ── Pruning results ──────────────────────────────────────────────
%
%   - Pruned architecture: [[3, 0], [3, 0], [2, 0]] in PyKAN's
%     (sum\_units, mult\_units) per-layer notation.
%   - No edges pruned: pre-prune edge analysis at threshold 0.01 finds
%     15 of 15 active (100\% survival). The model's edge-importance
%     distribution does not contain any below-threshold edges.
%   - Post-prune validation accuracy: 0.4667 (unchanged from Phase 2b
%     final, since no edges were removed).
%
% ── Symbolification rate: 100\% (15 of 15 edges) ───────────────────
%
%   - All 15 symbolified edges clear the 0.30 R\^2 threshold; zero
%     edges skipped during symbolic substitution. R\^2 distribution
%     across the 15 symbolified edges: min 0.734, median 0.990, max
%     1.000. Top-five edges all post R\^2 > 0.99 with the sin primitive
%     dominating the top of the ranking.
%
% ── Surviving features (3) ───────────────────────────────────────
%
%   - eth\_btc\_ratio (external / crypto-macro group, CoinMetrics-
%     sourced ETH/BTC alt-rotation signal). Selected in 57\% of CPCV
%     folds at the 28-fold level; most stable feature in the dataset
%     across all locked configurations to date.
%   - log\_returns\_lag6 (lag / autoregressive group, 6-day lag of log
%     returns). Selected in 54\% of CPCV folds. First time
%     log\_returns\_lag6 enters the symbolic representation; the
%     previous run had log\_returns\_lag1 as the lag-group entrant at
%     32\% selection frequency. The lag group has now contributed a
%     representative to two consecutive locked configurations (lag1
%     then lag6), validating the lag-features-in-MDA-pool design choice
%     as producing reproducible (if shifting) lag representation in the
%     symbolic formula.
%   - natgas\_ret\_30 (external / macro group, 30-day natural gas
%     return). Selected in 39\% of CPCV folds. Returning from the
%     previous run (where it was at 32\%). Natural gas is the only
%     macro-group feature to enter symbolic extraction in any of the
%     73-feature locked configurations.
%   - Group breakdown: 1 crypto-macro (eth\_btc\_ratio), 1 macro
%     (natgas\_ret\_30), 1 lag (log\_returns\_lag6), 0 TA, 0 on-chain,
%     0 mathematical (AFML Part 4). TA features dropped entirely from
%     the surviving set (was 2 of 5 in the previous run); the surviving
%     set is now distributed across three feature families rather than
%     four.
%
% ── Pre-symbolic vs. post-symbolic accuracy ──────────────────────
%
%   - Pre-symbolic (B-spline KAN, post-prune): 46.67\% (below 50\%
%     baseline).
%   - Post-symbolic (closed-form formula, after affine fine-tuning):
%     51.85\% (barely above 50\% baseline).
%   - Delta = +5.18 percentage points (post exceeds pre by 5.2 pp).
%   - Honest reading: both pre- and post-symbolic accuracies sit at or
%     near the 50\% baseline. The pre-symbolic 46.67\% is below chance:
%     the PyKAN model retrained on the deployment fold has not extracted
%     signal that beats the majority-class predictor (which would post
%     56.85\% by predicting class +1 always). The post-symbolic 51.85\%
%     is barely above chance and below both the majority baseline
%     (56.85\%) and the dynamic Phase-1 gate (56.93\%).
%   - The +5.18 pp gain is larger than the previous run's +1.5 pp gain
%     in magnitude, but the absolute level on which it operates
%     (sub-baseline pre-symbolic accuracy) means the gain represents a
%     recovery from below-chance to barely-above-chance rather than an
%     improvement from a meaningful baseline. The affine fine-tuning
%     step is doing genuine work (it consistently improves over the
%     pre-symbolic spline KAN), but the underlying model on this fold
%     has not found exploitable signal in the three input features.
%   - Cross-run trajectory of (pre-symbolic, post-symbolic, Delta):
%     (0.5564, 0.5263, -0.0301) -> (0.5564, 0.5940, +0.0376) -> (0.5940,
%     0.6090, +0.0150) -> (0.4667, 0.5185, +0.0518). The most recent
%     run has the largest gain magnitude but starts from the lowest
%     baseline.
%
% ── The decision function (numbered equation in the thesis) ──────
%
% Use a numbered equation environment to display the formula. Format
% reference (paste these terms into align* / cases for readability):
%
%   decision(x) =
%       + 3 (-284 eth\_btc\_ratio / 505
%            - (-1766 log\_returns\_lag6 / 179 - 1529/843)\^4 / 1000
%            - 434 tanh(3217 natgas\_ret\_30 / 944 + 2141/935) / 663
%            + 4511/913)\^2 / 283
%       - 351 sin( (797/768 - 3716 natgas\_ret\_30 / 745)\^3 / 109
%                  - 405 sin(4547 log\_returns\_lag6 / 949 + 147/754) / 508
%                  + 517 cos(2279 eth\_btc\_ratio / 686 + 4631/899) / 687
%                  + 433/74 ) / 877
%       + 197 cos( 10 (797/768 - 3716 natgas\_ret\_30 / 745)\^3 / 847
%                  - 625 sin(4547 log\_returns\_lag6 / 949 + 147/754) / 609
%                  + 31 cos(2279 eth\_btc\_ratio / 686 + 4631/899) / 32
%                  + 2734/385 ) / 828
%       + 455 tanh( 235 eth\_btc\_ratio / 924
%                   + 209 tanh(3217 natgas\_ret\_30 / 944 + 2141/935) / 706
%                   - 4610/959 ) / 932
%       + 11/19
%
%   P(up | x) = 1 / (1 + exp(-decision(x)))
%
% Coefficients are post-sympy.nsimplify(tolerance=1e-3) rationals. The
% structure has changed substantially from the previous configuration:
% instead of four outer terms each combining all five features through
% one of {sin, sin, tanh, tanh} wrappers, the new formula has four
% outer terms each using a different outer function: an outer squared
% term 3(...)\^2/283, an outer sin, an outer cos, and an outer tanh.
% The features are distributed unevenly: natgas\_ret\_30 appears in all
% four outer terms (via tanh and via cube primitives),
% log\_returns\_lag6 appears in three (the squared term and the two
% trigonometric terms), eth\_btc\_ratio appears in all four. The
% previous run's "each feature in all four outer terms" symmetry does
% not hold here.
%
% ── Functional-form pattern ──────────────────────────────────────
%
%   - The formula uses six distinct primitives: sin, cos, tanh, the
%     squared term (...)\^2, the cubed term (...)\^3, and the quartic
%     term (...)\^4. This is again substantively different from the
%     previous configuration's sin, tanh, atan, x\^3, |.| set:
%       * atan and |.| are gone. The previous run had atan attached to
%         log\_returns\_lag1 and natgas\_ret\_30 as a bounded smooth
%         saturator and |.| attached to natgas\_ret\_30 as a non-smooth
%         V-shape primitive. Neither appears in the v5 formula. The
%         non-differentiability flag from the previous configuration
%         is resolved: with |.| no longer in the primitive set, no
%         surviving feature posts an undefined gradient at the dataset
%         mean.
%       * cos returns. The previous run replaced cos with sin in every
%         outer-wrapper position; the v5 formula uses both, with cos
%         appearing in two positions (an outer cos and an inner cos
%         attached to eth\_btc\_ratio).
%       * x\^2 returns. Used as the outer wrapper of the first outer
%         term, with the inner argument combining eth\_btc\_ratio,
%         tanh(natgas\_ret\_30 + ...), and a quartic in
%         log\_returns\_lag6.
%       * x\^4 is new. The first time a quartic primitive survives the
%         R\^2 threshold in any locked configuration. Attached to
%         log\_returns\_lag6 in the first outer term.
%       * x\^3 survives, attached to natgas\_ret\_30 this time (was
%         attached to log\_returns\_lag1 in the previous run).
%       * tanh survives as both an outer wrapper (the fourth outer
%         term) and an inner primitive (tanh(3217 natgas\_ret\_30 /
%         944 + 2141/935) appears nested inside both the squared first
%         term and the fourth tanh term).
%   - The four outer wrappers (squared, sin, cos, tanh) are
%     functionally distinct. The KAN learned four different smooth
%     representations and combined them as a weighted sum. The previous
%     configuration's "two complementary internal representations x two
%     outer wrappers" symmetry does not hold in v5.
%
% ── Interpretive note (factual; economic interpretation in 5.3) ──
%
%   - The term-structure is no longer symmetric across features:
%     eth\_btc\_ratio appears in all four outer terms, natgas\_ret\_30
%     in all four (in different roles: tanh nested inside the squared
%     term, x\^3 inside the sin and cos terms, and tanh again as the
%     outer's inner argument), and log\_returns\_lag6 in three (the
%     squared term as a x\^4, and the two sine/cosine inner arguments).
%     The cross-run robustness of the "all features in all four outer
%     terms" pattern from the previous run does not survive the n=3
%     input dimensionality of v5.
%
% ── Numerical sensitivity at the dataset mean ────────────────────
%
% Partial derivatives evaluated at the mean of each feature, scaled by
% the feature's standard deviation, give the per-sigma effect on the
% decision function (logit). All three surviving features now have
% finite, well-defined gradients at the dataset mean (the previous
% run's natgas Abs-primitive non-differentiability is gone). Locked-run
% table (paste as a tabular environment in the thesis); values are TBD
% pending the locked-notebook sensitivity-cell rerun:
%
%   Feature           | mean | std | d/dx at mean | sigma-effect on logit | sigma Delta\_p (linearised)
%   eth\_btc\_ratio    | TBD  | TBD | TBD          | TBD                   | TBD
%   log\_returns\_lag6 | TBD  | TBD | TBD          | TBD                   | TBD
%   natgas\_ret\_30    | TBD  | TBD | TBD          | TBD                   | TBD
%
%   - Sensitivity table values are TBD pending the sensitivity
%     computation cell run in the locked notebook. The symbolic
%     decision function and its partial derivatives are reported in
%     the notebook output (Section 6.4 of main.ipynb); the
%     sigma-scaled effects need to be transcribed once the cell is
%     rerun with the current locked seed configuration. The
%     qualitative pattern from the symbolic form: natgas\_ret\_30
%     enters via tanh and x\^3 (both bounded primitives at the dataset
%     mean), log\_returns\_lag6 enters via x\^4 and the trigonometric
%     inner arguments (x\^4 contributes a steep gradient if the inner
%     argument is far from zero; at the dataset mean log\_returns\_lag6
%     ~ -0.0005 the x\^4 term contributes a small gradient),
%     eth\_btc\_ratio enters via the cos inner argument and the tanh
%     outer wrapper (both bounded and smooth at the dataset mean). The
%     locked formula is therefore everywhere-differentiable at the
%     dataset mean; the non-differentiability that previous-run readers
%     would have had to manage is resolved.
%   - Scale check. All three surviving features have interpretable
%     scales. eth\_btc\_ratio is bounded above by the historical
%     maximum (~0.13 in this sample); log\_returns\_lag6 is a 6-day log
%     return (bounded around +/- 0.30); natgas\_ret\_30 is a 30-day
%     return (typically +/- 0.20). No rescaling caveat is needed for
%     the formula.
%
% ── On-chain comment (Q2 evidence) ───────────────────────────────
%
%   - Zero on-chain features survive extraction. This finding is now
%     stable across the previous and current 73-feature locked
%     configurations. The Q2 answer hardens further: not only do
%     on-chain features fail to clear cross-fold stability (Section
%     4.5), they consistently fail to enter the locked symbolic
%     representation regardless of n\_top\_features configuration. The
%     earlier 62-feature locked-run observation that one on-chain
%     feature could carry the formula's strongest economic effect is
%     regime-specific and does not generalise to the expanded feature
%     universe.
%
% ── Lag-feature comment (Q2 evidence) ────────────────────────────
%
%   - A lag feature survives extraction for the second consecutive
%     locked configuration, but the specific lag has shifted
%     (log\_returns\_lag1 -> log\_returns\_lag6). The lag-features-in-
%     MDA-pool design choice (Section 3.4.4, advisor-driven, lag
%     features compete with engineered features rather than feeding
%     only into AR Logistic) admits lag features into the surviving
%     set in the 73-feature pool; the shifting lag identity across runs
%     suggests the autoregressive signal in BTC daily direction is not
%     concentrated at any single horizon but rotates across nearby lags
%     depending on the seed and feature-selection-strategy
%     configuration.
%
% ── Surviving-feature group breakdown ────────────────────────────
%
%   - 1 of 3 from crypto-macro (eth\_btc\_ratio).
%   - 1 of 3 from macro (natgas\_ret\_30).
%   - 1 of 3 from lag (log\_returns\_lag6).
%   - 0 of 3 from TA, on-chain, or mathematical (AFML Part 4) groups.
%
% ── Section closer ───────────────────────────────────────────────
%
%   - "The symbolic extraction pipeline produces a closed-form
%      decision function with 100\% symbolification rate (15 of 15
%      edges) on a three-feature subset (eth\_btc\_ratio,
%      log\_returns\_lag6, natgas\_ret\_30). Post-symbolic accuracy is
%      51.85\%, a 5.18 percentage-point gain over the pre-symbolic
%      46.67\%, but both numbers sit at or below the 50\% baseline:
%      the PyKAN model retrained on the deployment fold did not extract
%      signal that beats coin-flip. The formula uses six primitives
%      (sin, cos, tanh, x\^2, x\^3, x\^4) and four outer terms with
%      distinct outer wrappers (squared, sin, cos, tanh). The closed-
%      form approximation improves accuracy relative to the trained
%      spline network in absolute pp terms but the absolute level
%      remains sub-baseline, the Phase-1 dynamic gate fails (val\_acc
%      0.4815 < 0.5693), and the pre-prune gate also fails. The Abs(.)
%      non-differentiability flag from the previous configuration is
%      resolved: with |.| no longer in the primitive set, no surviving
%      feature posts an undefined gradient at the dataset mean. RQ4
%      conclusion: closed-form interpretability is achievable (100\%
%      symbolification rate, clean primitives, finite gradients
%      everywhere), but the locked-run model on the deployment fold
%      has not extracted signal that beats the majority baseline; the
%      symbolic representation faithfully captures a model that has
%      not found the signal."



% ==================================================================
% CHAPTER 5: DISCUSSION
% ==================================================================
\newpage
\section{Discussion}\label{sec:discussion}

% T-R7. Interpret. Link findings back to the mechanisms in Section 3.


\subsection{Interpreting the Results}\label{subsec:disc_results}

% Q1, Q3 evidence.
%
% Topics:
%   - The joint DSR / PBO finding. No model achieves DSR >= 0.95 (top
%     trained model is KAN at 0.2470; buy-and-hold posts a higher
%     median Sharpe of 2.2722 but is not included in the DSR
%     computation since it does not enter the selection pool). PBO
%     sits in the adversarial regime at 0.657, with 23 of 35 IS/OOS
%     partitions seeing the IS-winner underperform the OOS median.
%     The combination is informative: model selection in this
%     six-model cross-section is not just unreliable but actively
%     anti-informative; the IS-best is the OOS-loser in nearly two
%     of three partitions. The previous run's KAN strict-positive
%     bootstrap CI does NOT reproduce in the current configuration:
%     6 of 7 trained-model CIs cross zero, and only buy-and-hold's
%     CI (1.8903, 2.4685) stays strictly positive. KAN's CI moves
%     from the previous run's (0.0971, 1.2413) to a zero-crossing
%     (-0.4677, 1.1607) despite its median Sharpe rising from 0.4778
%     to 0.5588. The headline KAN-strict-positive finding from the
%     previous run is therefore not robust to the seed-count change.
%     The honest reading is "no trained model clears the multiple-
%     testing-corrected significance threshold; no trained-model
%     bootstrap CI stays strictly positive under uniform 5-seed
%     evaluation; model selection in this cross-section is
%     structurally adversarial; and a naive buy-and-hold strategy on
%     the same data posts a median Sharpe roughly 4x the best trained
%     model's figure". This is a stronger negative position than the
%     previous run could deliver: we have moved past both the moderate-
%     overfitting regime (PBO 0.34 -> 0.629 -> 0.657) AND the single-
%     model bootstrap-strict-positive evidence (which the v5 run
%     produced but v6 does not).
%   - PBO trajectory across the five locked configurations. PBO has
%     tracked 0.26 -> 0.37 -> 0.34 -> 0.629 -> 0.657 across the
%     62-feature, 73-feature 2-seed, 73-feature 3-seed-earlier,
%     73-feature 3-seed-previous, and current 73-feature 5-seed
%     runs. The first three configurations sit inside the moderate-
%     overfitting band [0.25, 0.40]; the last two are in the
%     adversarial regime (PBO > 0.5). The current configuration's
%     5-seed change does not move PBO back into the moderate band;
%     the adversarial regime is now stable across the v5 (3-seed)
%     and v6 (5-seed) runs, suggesting it is a real property of the
%     six-model cross-section under the 73-feature pool rather than
%     a single-run statistical artefact.
%   - Locating the variance: Random Forest is the PBO-anchor in v6.
%     The leave-one-out PBO (Section 4.4) shows that Random Forest's
%     removal leaves PBO unchanged at 0.657 (Delta +0.000), while
%     every other model's exclusion reduces PBO. Logistic Regression
%     is the largest destabiliser (Delta -0.286 to 0.371, back in
%     the moderate-overfitting band); AR Logistic and LSTM are joint
%     second largest (both Delta -0.200 to 0.457); KAN and XGBoost
%     are joint third (both Delta -0.114 to 0.543). The structural
%     finding is the Random Forest anchor: RF sits low enough in the
%     IS ranking across all 35 sub-path partitions that removing it
%     does not change the rotation among the remaining five models.
%     The previous run's "no anchor" pattern (every exclusion reduces
%     PBO) does not reproduce. Two readings across the LOO trajectory:
%     (a) the v5 no-anchor regime was the outlier across the locked
%     configurations evaluated; (b) the anchor identity varies with
%     the seed configuration (KAN-from-below in v4, no anchor in v5,
%     RF-from-below in v6), so the cross-section's IS/OOS rotation
%     structure is itself sensitive to within-model averaging
%     decisions.
%   - Consistency with EMH. The DSR < 0.95 finding aligns with semi-
%     strong-form EMH (Fama 1970) for BTC daily direction: under
%     leakage-free evaluation with multiple-testing correction, no
%     architecture in the comparison demonstrates a statistically
%     significant predictive edge. The fact that buy-and-hold posts a
%     higher median Sharpe (2.2722) than any trained model strengthens
%     this reading: in the 2015 to 2026 BTC sample, the unconditional
%     long position dominates every directional-trading strategy in
%     median Sharpe terms, suggesting that the asset's overall
%     trajectory carries information that direction-prediction
%     strategies discard by going neutral or short in volatile
%     regimes. The adversarial PBO finding hardens this further: not
%     only is no architecture significantly predictive, but the
%     cross-section's IS rankings are anti-informative about OOS
%     performance. This is the intellectually honest answer to Q1 and
%     consistent with \cite{chassot_audrino_2026}: properly fitted
%     baselines are hard to beat, and the literature's ML-superiority
%     results came from suboptimal fitting schemes that handicapped
%     the baselines. Cite: \cite{fama_1970}, \cite{chassot_audrino_2026}.
%   - The DeLong 2-of-15 finding is the new positive result. In the
%     v6 5-seed configuration, two pairwise AUC differences clear
%     alpha = 0.05: AR Logistic vs XGBoost (p = 0.0427) and Random
%     Forest vs XGBoost (p = 0.0318). The two near-significant pairs
%     at p ~ 0.054 in the previous 3-seed run clear the threshold
%     once the within-model averaging is over 5 seeds. Both
%     significant pairs include XGBoost as the lower-AUC half, with
%     XGBoost's pooled AUC (0.4942) sitting below 0.5. The structural
%     reading: XGBoost is significantly worse on classification AUC
%     than the two best-AUC trained models, but no positive AUC
%     difference above 0.5 has been demonstrated. The DeLong finding
%     does not propagate to a financial-side significance claim
%     (none of the significant AUC pairs has a positive bootstrap
%     CI on its corresponding path-Sharpe difference), but it is the
%     first statistically significant classification-side finding
%     across all five locked configurations and worth reporting as a
%     genuine positive result from the 5-seed methodology change.
%   - Conservative Predictions (Nabar and Shroff 2023): in low-signal
%     regimes, abstention beats prediction. The bet-sizing threshold
%     operationalises this: predictions with p ~ 0.50 produce bet
%     ~ 0, structural abstention without a separate model. The
%     locked-run mean P(Up) for KAN is 0.5421 against an empirical
%     base rate of 0.5685 (Delta = -0.0263, within the +/- 0.03
%     calibration tolerance), so the calibration audit does not flag
%     KAN. The bet-sizing curve is operating on well-calibrated
%     probabilities, and the negative DSR finding is therefore a
%     statement about signal strength, not about probability quality.
%     Only one model is flagged by the calibration audit in the
%     current configuration (XGBoost at Delta -0.0371, just over the
%     tolerance). The calibration footprint is comparable to the
%     previous run; despite the adversarial PBO regime and the loss
%     of the KAN strict-positive CI, calibration quality remains
%     well-controlled across the cross-section.
%   - Three-version pipeline evolution (advisor narrative): inflated
%     ROC-AUC -> degenerate all-Down predictions -> AFML. Two
%     specific failure modes motivated AFML adoption: uninformative
%     fixed-time-horizon labels and structural leakage that standard
%     splits cannot prevent. The locked-run results show that even
%     after AFML closes both leakage channels, the residual signal
%     is too weak for any compared model to clear DSR = 0.95, and
%     the IS rankings in the model cross-section are now anti-
%     informative under PBO; this is the AFML pay-off. The
%     methodology produces honest results (PBO 0.657 in the
%     adversarial regime) rather than the inflated metrics that
%     earlier (leaky) pipeline versions produced. Across the five
%     locked configurations evaluated to date, the PBO trajectory is
%     stable in the moderate band for the first three and adversarial
%     across the last two; the moderate-overfitting reading is now
%     understood as a small-sample artefact of the 2-seed and 3-seed
%     configurations, and the 5-seed v6 reading (PBO 0.657) is the
%     most-conservative locked figure to date.


\subsection{KAN Performance in Context}\label{subsec:disc_kan}

% Q3, Q4 evidence.
%
% Topics:
%   - KAN ranks 1st among trained models in median Sharpe (0.5588 in
%     v6, up from 0.4778 in v5), ahead of AR Logistic (0.3806), LSTM
%     (0.1541), XGBoost (0.0493), Random Forest (0.0287), and
%     Logistic Regression (-0.2173). Buy-and-hold's 2.2722 sits above
%     the entire trained-model cross-section. KAN's positioning is
%     stable across the two consecutive 73-feature locked
%     configurations: rank-1 trained Sharpe in both v5 (3-seed) and
%     v6 (5-seed) runs. KAN is also rank-1 on DSR (0.2470 in v6, up
%     from 0.2198 in v5) and posts the smallest median max-drawdown
%     among trained models (-0.2466 vs. -0.2172 for XGBoost in
%     second place). What does NOT survive the seed-count change is
%     the bootstrap-strict-positive CI: KAN's CI moves from (0.0971,
%     1.2413) in v5 to (-0.4677, 1.1607) in v6, now crossing zero
%     on the lower side. The headline KAN result from the previous
%     run is therefore not robust under 5-seed evaluation.
%   - Why the strict-positive CI does not reproduce. The mechanism is
%     a within-model variance change: averaging the per-(model,
%     split) calibrated probabilities over 5 seeds rather than 3
%     widens the bootstrap CI on the lower side more than it tightens
%     the median Sharpe. KAN's median Sharpe rises (0.4778 -> 0.5588)
%     and KAN's DSR rises (0.2198 -> 0.2470), but the path-Sharpe
%     std also rises (0.7592 -> 0.7539, essentially unchanged) and
%     the bootstrap resampling distribution on the 7 path-Sharpes
%     produces a CI lower tail that crosses zero. The honest reading:
%     KAN's positive median Sharpe is real (rank-1 trained, two
%     consecutive runs), but the bootstrap evidence on the path-level
%     that it is distinguishable from zero is sensitive to the
%     within-model averaging cardinality. A reader who weighted the
%     v5 strict-positive CI heavily should downweight that finding;
%     a reader who weighted the median Sharpe and DSR rankings should
%     note both continue to favour KAN.
%   - Classification-side metrics are no longer dominant. KAN posts
%     mean accuracy 0.5318 (mid-pack, ahead of XGBoost at 0.5185 and
%     LSTM at 0.5179 but behind AR Logistic at 0.5240 and Random
%     Forest at 0.5248), mean AUC 0.5103 (mid-pack, ahead of XGBoost
%     at 0.5008 but behind AR Logistic at 0.5204, Random Forest at
%     0.5166, and Logistic Regression at 0.5156), and median win
%     rate 0.5315 (third among trained models, behind AR Logistic
%     0.5318 and XGBoost 0.5122). The DeLong test (Section 4.3) now
%     finds 2 of 15 pairs significant at alpha = 0.05, but neither
%     involves KAN: the significant pairs are AR Logistic vs XGBoost
%     and Random Forest vs XGBoost. The previous configuration's
%     "KAN dominates on classification metrics" reading does not
%     survive into the current run; KAN is mid-pack across all three
%     classification-side metrics.
%   - Where KAN's rank-1 trained Sharpe comes from. Given that KAN
%     is mid-pack on AUC and mean accuracy, the rank-1 trained median
%     Sharpe must come from the bet-sizing translation rather than
%     from raw classification accuracy. Three diagnostics support
%     this reading: (a) KAN's mean P(Up) is 0.5421, within tolerance
%     of the 0.5685 empirical base rate at Delta = -0.0263, so the
%     bet-sizing curve operates on calibrated probabilities; (b)
%     KAN's median Sortino 0.3211 is second among trained models
%     (behind AR Logistic at 0.3894), indicating that KAN's
%     path-Sharpe distribution loads more variance on the upside than
%     the downside; (c) KAN's median max-drawdown -0.2466 is the
%     smallest among trained models, reflecting bet-sizing that
%     protects the strategy from the largest path-level moves. The
%     combination of well-calibrated probabilities, second-best
%     Sortino, and smallest max-drawdown produces the rank-1 trained
%     Sharpe and the rank-1 trained DSR, even though the bootstrap CI
%     no longer demonstrates strict-positive distinguishability from
%     zero.
%   - Tension to interpret: KAN is mid-pack on the LOO PBO table.
%     KAN's exclusion reduces PBO from 0.657 to 0.543 (Delta -0.114),
%     placing KAN tied with XGBoost as joint third-largest
%     destabiliser. The PBO-neutral model in v6 is Random Forest
%     (Delta +0.000); the v4 KAN-anchor finding does not reproduce
%     in either v5 (no anchor) or v6 (RF anchor). The reading: KAN
%     is the financially-best trained model in both consecutive
%     73-feature locked runs but is also part of the IS/OOS rotation
%     cluster contributing to the adversarial PBO baseline, not the
%     structural anchor stabilising the cross-section.
%   - Statistical-vs-economic gap. KAN's mean AUC of 0.5103 is mid-
%     pack in the cross-section, and the DeLong test against every
%     other model returns p > 0.09. The classification-side
%     superiority observed in earlier runs has not reproduced. What
%     does translate into a ranking position is the bet-sizing
%     transformation of KAN's well-calibrated probabilities: at a
%     base rate of 0.5685 and a KAN mean P(Up) of 0.5421, the
%     bet-sizing curve places conservatively-sized bets that protect
%     against the largest drawdowns. The honest reading: KAN's
%     classification advantage is not present in the current
%     configuration; its rank-1 Sharpe and DSR ranking comes from
%     the combination of well-calibrated bet sizing, smallest
%     max-drawdown, and second-best Sortino, not from a dominant
%     classification edge.
%   - Contrast with KASPER (Oad et al. 2025). KASPER reports R\^2 =
%     0.89 and Sharpe = 12.02 on individual stocks with regime
%     detection. This thesis posts median Sharpe = 0.56 for KAN on
%     BTC daily direction without a regime layer. The differences
%     are explanatory rather than competitive: KASPER uses regression
%     on individual stocks with regime-conditional architectures;
%     this thesis uses classification on a single asset (BTC) with
%     no regime layer and applies AFML's full statistical-correction
%     stack (DSR, PBO, DeLong) which KASPER does not. The roughly
%     21x Sharpe gap reflects asset class, target type, and
%     evaluation rigour rather than KAN architectural capacity.
%   - Contrast with VIX KAN (Cho et al. 2025). The VIX KAN paper
%     extracts symbolic formulas that reveal mean-reversion and
%     leverage effects, validated against domain knowledge. This
%     thesis's symbolic formula reveals a three-feature mixed-
%     primitive structure through sin, cos, tanh, x\^2, x\^3, and
%     x\^4 primitives, on a target (BTC daily direction) that is not
%     mean-reverting and where domain knowledge does not directly
%     validate any specific functional form. The symbolic-extraction
%     outputs are identical between the v5 (3-seed) and v6 (5-seed)
%     runs (PyKAN is retrained with a fixed internal seed on the
%     same training fold), so the symbolic finding is invariant
%     under the outer-loop seed change. The symbolic-extraction
%     machinery transfers across asset classes; the interpretive
%     yield depends on whether the underlying target has a structure
%     that domain knowledge can corroborate.


\subsection{Symbolic Extraction as Contribution}\label{subsec:disc_symbolic}

% Q4 evidence.
%
% Invariance note (v6, 5 seeds). The symbolic-extraction outputs are
% identical to the previous 3-seed locked run. PyKAN is retrained
% from scratch on the same training fold (split 27) using a fixed
% internal seed, so the closed-form formula, the architecture
% [3, 3, 2], the 100\% symbolification rate, the 46.67\% / 51.85\%
% pre/post symbolic accuracies, and the phase-by-phase validation
% accuracies are all invariant under the outer-loop seed-count
% change. All claims in this subsection therefore refer to a
% symbolic representation that is unchanged across the v5 and v6
% configurations.
%
% Topics:
%   - Interpretability is a separable contribution from predictive
%     accuracy. The locked-run formula has 100\% symbolification rate
%     (15 of 15 edges, no skips) and a closed-form expression on three
%     features. Its post-symbolic accuracy of 51.85\% is a 5.18
%     percentage-point gain over the spline KAN's pre-symbolic 46.67\%,
%     but both numbers sit at or below the 50\% baseline and well
%     below the 56.85\% majority baseline. The closed-form decision
%     function provides three artefacts a black-box model cannot:
%     per-feature symbolic derivatives (Section 4.6 numerical
%     sensitivity), term-structure decomposition (each feature appears
%     in at least three of four outer terms), and an audit trail from
%     input to probability that a domain expert can inspect. The fact
%     that the symbolic form outperforms the spline form is consistent
%     across all five locked configurations evaluated to date (-3.0
%     pp, +3.8 pp, +1.5 pp, +5.2 pp, +5.2 pp; the v6 run reproduces
%     the v5 figure exactly since PyKAN is retrained from scratch
%     with a fixed internal seed on the same training fold and is
%     invariant under the outer-loop seed-count change); the symbolic substitution
%     appears to act as a smoothing regulariser on the per-edge
%     B-spline shapes, and the affine fine-tuning step then tunes the
%     per-edge scale and bias to the validation distribution. What
%     does not survive in the current run is the absolute level: post-
%     symbolic 51.85\% sits below the majority baseline 56.85\%,
%     meaning the closed-form formula does not beat the always-predict-
%     class-+1 predictor on this fold.
%   - The methodology change (n\_top\_features=3). The locked notebook
%     call uses n\_top\_features=3 (down from 5 in the previous run),
%     restricting the symbolic extraction input to the three most-
%     stable features by cross-fold MDA frequency: eth\_btc\_ratio
%     (57\%), log\_returns\_lag6 (54\%), natgas\_ret\_30 (39\%). The
%     motivation for the change (as inferable from the codebase:
%     n\_top\_features is exposed as a parameter to
%     run\_symbolic\_extraction and the
%     feature\_selection\_strategy="stability" default ranks by 28-fold
%     selection frequency) is to tighten the formula to a humanly-
%     auditable size on the deployment fold; the trade-off is that the
%     input dimensionality drop forces a narrower [3, 3, 2]
%     architecture and reduces the data PyKAN has to find signal on.
%     The current run is the first locked configuration where the
%     methodology change is made explicit; the previous three runs
%     used the default n\_top\_features=5. The honest reading: the
%     smaller input set produces a cleaner formula (100\% sym rate,
%     six clean primitives, no |.| non-differentiability) but on this
%     fold also produces a model that does not beat the majority
%     baseline.
%   - The post-symbolic accuracy gain in magnitude vs. level. The
%     +5.18 pp gain in the current configuration is the largest of
%     any locked run in absolute pp terms, but operates from the
%     lowest baseline (pre-symbolic 46.67\% vs. previous-run 59.40\%,
%     55.64\%, 55.64\% for the three prior runs). The reading: the
%     affine fine-tuning step is doing genuine work on the deployment
%     fold, recovering above-chance behaviour from a pre-symbolic
%     representation that was below chance. But the magnitude of the
%     recovery does not get the formula above the majority baseline,
%     so the symbolic representation faithfully captures a model that
%     has not found the signal. Per-fold extraction (L3 in 5.4) would
%     let the analysis distinguish between fold-specific gains and a
%     robust symbolic-smoothing effect; this remains the highest-
%     leverage methodological extension available to the project.
%   - The Phase-1 dynamic gate failure. The Phase 1 (Adam) validation
%     accuracy of 0.4815 falls below the dynamic gate max(0.53,
%     base\_rate + 0.01) = 0.5693 by 8.78 percentage points. The
%     dynamic gate is a methodology improvement over the previous
%     fixed 0.53 threshold: it adjusts to the actual class balance in
%     the validation fold rather than assuming the dataset-level base
%     rate, which makes it sensitive to per-fold imbalance. The pre-
%     prune gate (val\_acc < 0.5359, where 0.5359 = 0.5259 majority
%     baseline + 0.01 margin) also fails at val\_acc 0.4667. Both
%     gates fail in the current configuration, the first locked run
%     where the gate-failure is structural rather than recoverable.
%     The previous run's gate failure (Phase 1) was followed by LBFGS-
%     phase recovery to 0.5940; the current run does not recover above
%     the 50\% baseline at any of the four training phases. The
%     reading: the symbolic-extraction pipeline can still deliver a
%     usable closed-form formula on a model that did not pass its own
%     gates (the structural transparency of the decision function does
%     not depend on classification accuracy), but the gate failures
%     are honest signals that the underlying PyKAN model has not
%     extracted directional signal on this fold.
%   - Mixed-primitive functional form. The decision function uses six
%     distinct primitives (sin, cos, tanh, x\^2, x\^3, and x\^4),
%     arranged in four outer terms with structurally distinct outer
%     wrappers (squared, sin, cos, tanh). The previous run used sin,
%     tanh, atan, x\^3, and |.|; the current run replaces atan and
%     |.| (which had introduced the natgas non-differentiability
%     flag) with cos, x\^2, and x\^4. The new primitive set is
%     entirely differentiable; the previous run's non-differentiability
%     flag is resolved. The x\^4 primitive is new across all four
%     locked configurations and attaches to log\_returns\_lag6 in the
%     squared outer term. The four-outer-term structure recurs from
%     the previous run, suggesting PyKAN's pruning step produces a
%     four-outer-term representation under both the [5, 2, 2] and
%     [3, 3, 2] architectures, but the symmetry across outer terms is
%     lost in the current configuration: each outer wrapper is a
%     different function (squared, sin, cos, tanh) rather than two
%     pairs (sin x 2, tanh x 2) as in the previous run.
%   - Cross-run primitive-set instability. The primitive set has now
%     changed in every pair of consecutive locked configurations:
%     62-feature \{sin, cos, tanh, x\^2, x\^3\} -> 73-feature 2-seed
%     \{sin, tanh, atan, x\^3, |.|\} -> 73-feature 3-seed previous
%     \{sin, tanh, atan, x\^3, |.|\} -> current \{sin, cos, tanh, x\^2,
%     x\^3, x\^4\}. The cross-run robust subset is \{sin, tanh, x\^3\};
%     the cross-run rotating subset is \{cos, atan, |.|, x\^2, x\^4\}.
%     The reading: the brute-force fitting library tries roughly 14
%     candidates per edge and keeps the highest-R\^2 match per fold,
%     so different feature sets and folds land on substantively
%     different primitive combinations. Reporting this cross-run
%     instability is itself a finding about the symbolic-extraction
%     methodology rather than about the underlying BTC direction
%     process: the formula's functional signature is regime- and
%     configuration-conditional, not stable.
%   - Numerical sensitivity at the dataset mean. All three surviving
%     features have finite, well-defined gradients at the dataset
%     mean (the previous run's natgas Abs-primitive non-
%     differentiability is gone). The sensitivity table values are
%     TBD pending the locked-notebook sensitivity-cell rerun (see
%     Section 4.6 numerical sensitivity bullet). The qualitative
%     pattern from the symbolic form: natgas\_ret\_30 enters via tanh
%     and x\^3 (both bounded at the dataset mean), log\_returns\_lag6
%     enters via x\^4 and the trigonometric inner arguments (x\^4
%     contributes a steep gradient if the inner argument is far from
%     zero; at the dataset mean log\_returns\_lag6 ~ -0.0005 the x\^4
%     term contributes a small gradient), eth\_btc\_ratio enters via
%     the cos inner argument and the tanh outer wrapper (both bounded
%     and smooth at the dataset mean). The locked formula is therefore
%     everywhere-differentiable at the dataset mean; the non-
%     differentiability that previous-run readers would have had to
%     manage is resolved.
%   - The on-chain finding hardens further across configurations. No
%     on-chain feature survives extraction in any of the three
%     73-feature locked configurations. The previous 62-feature
%     locked run retained exchange\_supply\_pct as the single on-chain
%     feature in the surviving four; all three subsequent runs at 73
%     features have produced zero on-chain survivors. The Q2 answer
%     therefore reads: on-chain features may carry signal in some
%     narrowly-defined configurations (the 62-feature pool with
%     exchange\_supply\_pct at 46\% selection frequency), but the
%     broader configuration with 20 macro features and 10 lag
%     features competing in the MDA pool produces zero on-chain
%     survivors at any reasonable seeds configuration and at any
%     n\_top\_features setting. The on-chain free lunch fails at the
%     cross-fold stability level (no on-chain feature reaches 50\%
%     selection frequency across the 28 folds in any run) and at the
%     deployment-fold level (no on-chain feature reaches the symbolic
%     decision function on split 27 in any 73-feature run). The
%     earlier locked-run observation that one on-chain feature could
%     carry the formula's strongest effect was regime-specific and
%     does not generalise.
%   - The lag-feature finding now reproduces across two configurations.
%     A lag feature survives into the symbolic representation in both
%     the previous run (log\_returns\_lag1 at 32\% selection
%     frequency) and the current run (log\_returns\_lag6 at 54\%
%     selection frequency, moderate stability). The specific lag has
%     rotated (lag1 -> lag6) but the lag group's presence in the
%     surviving set is now reproducible. The lag-features-in-MDA-pool
%     design choice (Section 3.4.4, advisor-driven, lag features
%     compete with engineered features rather than feeding only into
%     AR Logistic) admits lag features into the surviving set across
%     configurations; the rotation of the specific lag identity
%     suggests the autoregressive signal in BTC daily direction is not
%     concentrated at any single horizon but rotates across nearby
%     lags depending on the feature-selection-strategy and seed
%     configuration. This is the cleanest reproducible Q2 finding so
%     far: lag features are consistently part of the surviving set in
%     the expanded 73-feature universe, even though no specific lag
%     is consistently the surviving lag.
%   - Contrast with VIX KAN (Cho et al. 2025). The VIX KAN paper
%     extracts symbolic formulas that reveal mean-reversion and
%     leverage effects, validated against domain knowledge. This
%     thesis's locked formula reveals a three-feature mixed-primitive
%     structure on a target (BTC daily direction) that is not mean-
%     reverting and where domain knowledge does not directly validate
%     any specific functional form. The cross-run variation in the
%     dominant feature and its sign is itself an honest finding about
%     the symbolic-extraction methodology: the formula's economic
%     interpretation is fold- and configuration-conditional, and the
%     current run's formula is faithful to a model that did not
%     extract directional signal on the deployment fold.
%   - Limitations specific to symbolic extraction (preview to 5.4):
%       * Fold-specific. Different folds may produce different
%         formulas with different surviving features; per-fold
%         extraction is L3 in 5.4. The cross-run change from
%         (eth\_btc\_ratio, bb\_width, skewness, log\_returns\_lag1,
%         natgas\_ret\_30) to (eth\_btc\_ratio, log\_returns\_lag6,
%         natgas\_ret\_30) is direct evidence that fold- and
%         configuration-specificity is a real limitation rather than
%         a theoretical concern.
%       * Small sample. 540 training observations on split 27 is not
%         large; PyKAN's per-edge R\^2 fits stabilise on this size
%         but recover signal that the spline KAN already captured
%         rather than discovering new structure. With only 3 input
%         features (down from 5), the input dimensionality is now low
%         enough that some signal in the dropped features may have
%         been informative.
%       * R\^2 threshold sensitivity. The 0.30 threshold (lowered
%         from 0.50) admits the lowest-R\^2 edge at 0.734 in the
%         current run, which is much cleaner than the previous run's
%         0.5746 minimum. The 100\% symbolification rate is itself a
%         quality improvement relative to the previous run's 93\%, and
%         the R\^2 distribution is tighter at the bottom end.
%       * The dynamic majority-baseline gate is a methodology
%         improvement that now fails in both phases. The previous
%         fixed-0.53 gate failed at val\_acc 0.5038 in the previous
%         run but the model recovered above baseline in subsequent
%         phases. The new dynamic gate (max(0.53, base\_rate + 0.01) =
%         0.5693 in this run) fails by a much wider margin (val\_acc
%         0.4815), and the pre-prune gate (0.5359) also fails. The
%         diagnostic improvement (dynamic threshold) has surfaced a
%         real problem (sub-baseline learning on this fold) that the
%         previous fixed threshold would have masked partially.
%       * The natgas\_ret\_30 non-differentiability flag is resolved.
%         The current run's primitive set (sin, cos, tanh, x\^2,
%         x\^3, x\^4) is entirely differentiable. The previous run's
%         |.| primitive that introduced the SymPy re/im decomposition
%         issue at the dataset mean is no longer in the surviving
%         formula. Sensitivity at the dataset mean is now finite for
%         every surviving feature.
%       * Architectural parity has been resolved. The single-hidden-
%         layer / width1 <= 6 configuration matches the CPCV-
%         benchmarked KAN, so the formula describes the architecture
%         the benchmark numbers describe, not a simplified surrogate.
%         The current run's [3, 3, 2] architecture uses width1 = 3
%         (down from 5) and width2 = 3 (up from 2), producing 15
%         active edges and a four-outer-term decision function that
%         fits one printed page in the thesis.


\subsection{Limitations}\label{subsec:disc_limitations}

% Six honest limitations, from defence L1 to L5 plus new L6.
%
% Topics:
%   - L1. Computational power. Currently 5 seeds across all 6 models
%     (uniform, producing 840 = 28 x 6 x 5 prediction entries; up
%     from 3 seeds / 504 entries in the previous run), 30 Optuna
%     trials per tuned model per fold, KAN width capped at 6. With
%     more compute: more seeds, more trials, wider HP ranges.
%     The N=8 / 28-split configuration alone roughly doubles the
%     per-experiment runtime relative to N=6 / 15 splits, and the
%     per-fold MDA inside Optuna's TPE inner loop adds roughly 40\%
%     on top.
%   - L2. Daily OHLCV only. Discards intraday microstructure. Pipeline
%     is timeframe-agnostic; hourly extension would multiply event
%     count and unlock microstructure features.
%   - L3. Symbolic extraction is fold-specific and configuration-
%     specific. The locked formula is the result of running the
%     extraction on split 27 only with n\_top\_features=3. Per-fold
%     extraction (28 formulas, one per CPCV split) would let the
%     analysis count which features and primitives recur (structural
%     signal vs. one-off noise). The cross-run variation observed to
%     date is substantial: surviving feature sets have rotated
%     \{exchange\_supply\_pct, atr, sadf, oil\_ret\_30\} ->
%     \{eth\_btc\_ratio, bb\_width, skewness, log\_returns\_lag1,
%     natgas\_ret\_30\} -> \{eth\_btc\_ratio, log\_returns\_lag6,
%     natgas\_ret\_30\} across the three feature-set configurations,
%     with only eth\_btc\_ratio recurring across all three. The
%     primitive set has rotated even more (see Section 5.3). Per-fold
%     extraction is the highest-leverage methodological extension
%     available to the project: it would let the analysis report "X
%     features recurred in N of 28 folds" instead of "X features
%     recurred in the deployment fold only", which is structurally a
%     much stronger claim.
%   - L4. KAN width cap is binding for interpretability, not capacity.
%     The width1 <= 6 cap is set so the extracted symbolic formula
%     stays humanly readable. The current run's pruned architecture
%     [[3, 0], [3, 0], [2, 0]] uses width1 = 3 and width2 = 3,
%     producing 15 active edges and a four-outer-term decision
%     function that fits one printed page in the thesis. The previous
%     run's architecture [[5, 0], [2, 0], [2, 0]] used width1 = 5 and
%     width2 = 2 with 14 edges and a different four-outer-term
%     structure. A larger width1 cap would let the model fit more
%     complex patterns but would also produce a formula too long to
%     interpret. This is an explicit interpretability-vs-capacity
%     trade-off, not a generic constraint imposed by sample size or
%     compute. The recurrence of the four-outer-term structure across
%     configurations suggests the cap is not binding for the discovery
%     of the structure itself, only for the readability of the
%     resulting formula.
%   - L5. No regime-conditional analysis. BTC has had approximately 5
%     distinct regimes over 2014 to 2026 (the 2017 bubble, 2018-2019
%     bear, 2020 COVID crash, 2021 bull, 2022 FTX-era contraction,
%     2023-2026 recovery). Full-sample Sharpe averages over them. A
%     regime-conditional analysis is left for future work, especially
%     given the joint reading "DSR < 0.95 + PBO 0.657 adversarial":
%     full-sample non-significance combined with adversarial cross-
%     section ranking does not rule out within-regime predictability,
%     and the buy-and-hold comparison (median Sharpe 2.27 with max
%     drawdown -0.9998) makes the regime-conditional question
%     financially salient even when full-sample-significance is
%     unattainable.
%   - L6. Sub-baseline accuracy on the deployment fold. The locked-run
%     symbolic extraction posts pre-symbolic accuracy 46.67\% (below
%     50\% baseline) and post-symbolic accuracy 51.85\% (just above
%     50\% baseline but below the 56.85\% majority baseline). Both
%     the Phase-1 dynamic gate (0.5693) and the pre-prune gate
%     (0.5359) fail on the deployment fold. The honest reading: the
%     symbolic representation faithfully captures a PyKAN model that
%     has not extracted directional signal on this fold. This is a
%     real limitation: the symbolic-extraction contribution is
%     methodological (closed-form, 100\% sym rate, finite gradients
%     everywhere, six clean primitives), not predictive on the
%     deployment fold itself. A future run could (a) revert
%     n\_top\_features to 5 to give PyKAN more data dimensionality on
%     the deployment fold, (b) test alternative fold selections to
%     find one where PyKAN clears its gates, or (c) accept the sub-
%     baseline result and report it as evidence that closed-form
%     extraction works even on weakly-learning models. The previous
%     locked configuration's |.| primitive non-differentiability flag
%     is now resolved; the current limitation list trades the non-
%     differentiability flag for the sub-baseline accuracy flag.



% ==================================================================
% CHAPTER 6: CONCLUSION
% ==================================================================
\newpage
\section{Conclusion}\label{sec:conclusion}

% Cochrane: do not restate all findings. T-R16: end with resonance, not
% a hedge. Target <= 10\% of the textual part (~3.5 pages).
%
% Suggested four-paragraph structure:
%
%   - One paragraph stating the contribution: this thesis applies the
%     complete AFML pipeline (CUSUM, TBL, FFD, sample weights, CPCV,
%     DSR, PBO, DeLong) to BTC daily direction prediction on a
%     73-feature universe spanning technical, mathematical (AFML Part
%     4), macro, crypto-macro, on-chain, and autoregressive-lag
%     families. It benchmarks KAN against five models (AR Logistic,
%     Logistic Regression, Random Forest, XGBoost, LSTM) under
%     identical CPCV splits with N=8 / k=2 / 28 splits / 7 backtest
%     paths and 5 seeds across all six models (840 prediction entries
%     from 28 splits x 6 models x 5 seeds; the v6 locked run; the
%     previous v5 run used 3 seeds / 504 entries), and extracts the
%     first
%     closed-form classification formula from a CPCV-trained KAN under
%     the full AFML evaluation. The current locked configuration adds
%     two methodology improvements relative to earlier runs: a dynamic
%     majority-baseline gate (max(0.53, base\_rate + 0.01)) that
%     replaces the fixed Phase-1 53\% threshold, and an exposed
%     n\_top\_features parameter on the symbolic-extraction call so
%     the input dimensionality of the closed-form formula is a
%     configurable choice rather than a hardcoded default. The
%     symbolic formula is the novel deliverable; the six-model
%     benchmark is the methodological scaffolding that validates it.
%   - One paragraph stating the headline: under leakage-free evaluation
%     with multiple-testing correction, the literature's 85\% to 95\%
%     accuracy claims do not survive. Top DSR among trained models is
%     KAN at 0.2470, far below the 0.95 significance threshold. PBO at
%     0.657 places model selection in the adversarial regime (PBO >
%     0.5): 23 of 35 IS/OOS partitions see the IS-winner underperform
%     the OOS median, and the leave-one-out PBO analysis shows that
%     Random Forest is the PBO-anchor (its exclusion leaves PBO
%     unchanged at 0.657), while every other model's exclusion reduces
%     PBO. Logistic Regression is the largest destabiliser
%     (Delta -0.286 to 0.371). The cross-section's IS rankings are
%     anti-informative about OOS performance. Under uniform 5-seed
%     evaluation, 6 of 7 bootstrap 95\% CIs cross zero; only
%     buy-and-hold's CI (1.8903, 2.4685) stays strictly positive. The
%     previous run's KAN strict-positive bootstrap CI does NOT
%     reproduce in the 5-seed configuration: KAN's CI moves from
%     (0.0971, 1.2413) to (-0.4677, 1.1607). Buy-and-hold posts a
%     higher median Sharpe (2.2722) than every trained model with the
%     cost of a -99.98\% maximum drawdown, illustrating how punishing
%     the AFML correction layer is for trading strategies that try to
%     add tactical direction on top of an underlying drift. The honest
%     answer to Q1: no architecture demonstrates a statistically
%     significant predictive edge after correcting for selection bias;
%     KAN posts the strongest single-model evidence (rank-1 trained
%     Sharpe, rank-1 trained DSR across two consecutive runs), but
%     its bootstrap CI now crosses zero. The DeLong test reaches
%     2-of-15 significant pairs in the 5-seed configuration (AR
%     Logistic vs XGBoost and Random Forest vs XGBoost, both at
%     alpha=0.05, both with XGBoost as the lower-AUC half), the first
%     classification-side significance in any locked configuration.
%     The PBO trajectory across the five locked configurations
%     evaluated to date (0.26 -> 0.37 -> 0.34 -> 0.629 -> 0.657) was
%     stable in the moderate-overfitting band for the first three
%     runs and adversarial across the last two; the v6 5-seed run
%     produces the strongest negative cross-section evidence to date
%     (adversarial PBO baseline, zero trained-model strict-positive
%     CIs) alongside the first DeLong-significant classification
%     finding.
%   - One paragraph on what does survive: the symbolic extraction
%     pipeline produces a humanly auditable closed-form decision
%     function on three features (eth\_btc\_ratio, log\_returns\_lag6,
%     natgas\_ret\_30) with 100\% symbolification rate (15 of 15
%     edges, no skips) and post-symbolic validation accuracy 51.85\%,
%     a 5.18 percentage-point gain over the spline KAN's 46.67\%. The
%     formula uses six primitives (sin, cos, tanh, x\^2, x\^3, x\^4)
%     and four outer terms with structurally distinct outer wrappers
%     (squared, sin, cos, tanh). The previous run's |.|-primitive
%     non-differentiability flag is resolved: all three surviving
%     features have finite, well-defined gradients at the dataset
%     mean. The lag-feature finding now reproduces across two
%     consecutive 73-feature locked configurations (log\_returns\_lag1
%     then log\_returns\_lag6), validating the lag-features-in-MDA-
%     pool design choice; the specific lag rotates but the lag group's
%     presence in the surviving set does not. No on-chain feature
%     survives extraction in any of the three 73-feature locked
%     configurations evaluated: the Q2 "on-chain free lunch" hypothesis
%     fails at both the cross-fold stability level (no on-chain
%     feature reaches 50\% selection frequency across the 28 folds in
%     any run) and the deployment-fold level (no on-chain feature
%     reaches the symbolic decision function on split 27 in any
%     73-feature run). The closed-form representation does not cost
%     validation accuracy on this run (in fact gains 5.2 pp; the
%     symbolic-extraction outputs are identical to the previous
%     3-seed run because PyKAN is retrained from scratch with a
%     fixed internal seed on the same training fold, so the
%     symbolic finding is invariant under the outer-loop seed-count
%     change)
%     but the absolute level remains below the majority baseline
%     (51.85\% < 56.85\%), meaning the symbolic representation
%     faithfully captures a model that did not extract directional
%     signal on the deployment fold. Interpretability is a separable
%     contribution from predictive accuracy and does not need to wait
%     for predictive significance to be valuable: 100\% symbolification
%     rate, six clean primitives, finite gradients everywhere, and a
%     closed-form formula on three economically interpretable features
%     are themselves the contribution.
%   - One closing sentence with resonance. T-R16 example template:
%     "X is not a constraint, but a catalyst, for Y." Adapt for this
%     thesis: "Methodological honesty is not a constraint on financial
%     machine learning; it is the prerequisite for treating
%     interpretability as a contribution in its own right, and for
%     reading buy-and-hold's dominance over every trained model as a
%     result rather than a failure."



% ==================================================================
% CHAPTER 7: FUTURE WORK
% ==================================================================
\newpage
\section{Future Work}\label{sec:future}

% Cochrane: do not write your grant application here. Concrete,
% actionable directions only.
%
% Topics:
%   - MultKAN (KAN 2.0) on the last fold. Multiplication nodes enable
%     discovery of multiplicative interactions. Same symbolic pipeline.
%     Cite: \cite{liu_kan2_2024}.
%   - Higher-frequency data. Hourly bars would multiply CUSUM events
%     approximately 24 times and unlock microstructure features.
%     Pipeline timeframe-agnostic.
%   - Per-fold symbolic extraction. Count which features and primitives
%     recur across all 28 folds to distinguish structural signal from
%     one-off noise.
%   - Regime-conditional analysis. Decompose Sharpe by regime
%     (approximately 5 distinct regimes 2014 to 2026).
%   - Meta-labeling layer. A second classifier on the best primary
%     classifier's output, trained to predict whether to take the trade
%     (AFML Ch. 3). Singh and Joubert (2019) provide empirical evidence
%     that meta-labeling improves signal efficacy across asset classes;
%     applying it on top of this thesis's calibrated probabilities is
%     the natural next layer. Cite: \cite{singh_joubert_2019}.
%   - Alternative assets. ETH, gold for symbolic-formula comparison
%     (different on-chain availability).
%   - Walk-forward vs. CPCV. Compare conclusions under both protocols.



% =============================================
% BIBLIOGRAPHY
% =============================================
\newpage
\phantomsection
\thispagestyle{plain}
\addcontentsline{toc}{section}{Bibliography}
\bibliographystyle{agsm}
\nocite{*}
\bibliography{mfw_references}


% =============================================
% APPENDICES
% =============================================
\newpage
\phantomsection
\thispagestyle{plain}

\appendix
\section{Appendices}\label{sec:appendices}

\subsection{Pipeline Architecture Diagram}\label{app:pipeline}
% Full pipeline flowchart: data -> labelling -> weights -> features
% -> alignment -> CPCV -> preprocessing -> models -> evaluation
% -> symbolic.

\subsection{CPCV Split Details}\label{app:cpcv_splits}
% Group boundaries, train/test timelines for all 28 splits, leakage
% audit.

\subsection{Hyperparameter Search Spaces}\label{app:hyperparams}
% Optuna search spaces for all tuned models.

\subsection{Per-Split Classification Reports}\label{app:split_reports}
% 28 splits x 6 models.

\subsection{Full Feature List}\label{app:features}
% Complete table of 73 features with descriptions, sources, parameters.

\subsection{Symbolic Extraction Detailed Output}\label{app:symbolic}
% Full PyKAN training logs, edge-by-edge R^2, pruned diagrams,
% unsimplified formulas.

\subsection{Leakage Prevention Summary}\label{app:leakage_summary}
% T-R15 / T-R17 (write for reviewers; easy to teach). Moved to the
% appendix during the page-budget pass; lives here as a single-page
% reference for defenders and reviewers who want a per-channel audit.
%
% Table suggestion: leakage prevention summary
%   Risk                                  | Mitigation                                                           | Where
%   FFD d* on full data                   | per-fold training-only estimation                                    | preprocessing.py
%   Scaler on full data                   | RobustScaler on training fold only                                   | preprocessing.py
%   Feature selection on full data        | multi-model MDA on training fold only                                | preprocessing.py
%   HP tuning on full data                | nested Optuna inside training fold (3-fold purged)                   | tuning.py
%   Overlapping TBL labels                | three-condition purging                                              | cv.py
%   Serial correlation across CV boundary | embargo after test boundary                                          | cv.py
%   Calibration on test data              | held-out 20\% of training fold, never test                           | calibration.py
%   XGB early-stop + cal shared cell      | acknowledged; only ensemble size affected                            | pipeline.py
%   On-chain look-ahead                   | CoinMetrics shifted by 1 day                                         | external\_features.py
%   CUSUM threshold on full data          | minor approximation; acknowledged, negligible impact                 | labeling.py
%   NaN imputation across train-test      | ffill().bfill() independent within each partition; FFD columns drop  | preprocessing.py


% =============================================
% End Document
% =============================================
\end{document}